%% Based on a TeXnicCenter-Template by Tino Weinkauf.
%%%%%%%%%%%%%%%%%%%%%%%%%%%%%%%%%%%%%%%%%%%%%%%%%%%%%%%%%%%%%

%%%%%%%%%%%%%%%%%%%%%%%%%%%%%%%%%%%%%%%%%%%%%%%%%%%%%%%%%%%%%
%% HEADER
%%%%%%%%%%%%%%%%%%%%%%%%%%%%%%%%%%%%%%%%%%%%%%%%%%%%%%%%%%%%%
\documentclass[oneside,11pt,review,usenames,dvipsnames]{report}
% Alternative Options:
%	Paper Size: a4paper / a5paper / b5paper / letterpaper / legalpaper / executivepaper
% Duplex: oneside / twoside
% Base Font Size: 10pt / 11pt / 12pt
\setlength{\parskip}{\baselineskip}

%% Packages
\usepackage[letterpaper, paperwidth=170mm, paperheight=225mm, bottom=20mm, right=20mm]{geometry}
\usepackage{amsmath}
\usepackage{amssymb}
\usepackage{amsthm} %paquete para s\'imbolo matem\'aticos
\usepackage[spanish,es-nodecimaldot]{babel} %francais, polish, spanish, ...
\usepackage[utf8]{inputenc}
\usepackage{enumerate}
\usepackage{fancyhdr} %%Fancy headings
\usepackage[T1]{fontenc}
\usepackage{graphicx}
\usepackage{bm}
%\usepackage{subfig} %%Subfigures inside a figure
%\usepackage{pst-all} %%PSTricks - not useable with pdfLaTeX
\usepackage[pdftex,
            pdfauthor={CARLOS OMAR PARDO GOMEZ},
            pdftitle={UN MODELO BAYESIANO Y NO PARAMETRICO DE REGRESION SOBRE CUANTILES},
            pdfsubject={TESIS PARA EL TITULO DE LICENCIADO},
            pdfkeywords={BAYESIANO CUANTILES NO-PARAMETRICO MODELO ESTADISTICA REGRESION},
            pdfproducer={Latex con hyperref},
            pdfcreator={pdflatex}]{hyperref}
\usepackage[ansinew]{inputenc}
\usepackage{lmodern} %Type1-font for non-english texts and characters
\usepackage{longtable} %%For tables, that exceed one page
\usepackage[round]{natbib}
\usepackage{subfig} %para poner subfiguras
\usepackage[nottoc]{tocbibind}
\usepackage{url}
\usepackage{fixltx2e}
\usepackage{dsfont}
\usepackage{fontawesome}
\usepackage{float}
\usepackage{caption}
\usepackage[spanish,onelanguage,ruled]{algorithm2e}

\usepackage{amsmath}
\usepackage{amssymb}
\usepackage{amsthm}
\usepackage{mathtools}

\makeatletter
\def\thm@space@setup{%
  \thm@preskip=20pt \thm@postskip=20pt
}
\renewcommand{\@algocf@capt@plain}{above}% formerly {bottom}
\renewcommand\fs@ruled{\def\@fs@cfont{\bfseries}\let\@fs@capt\floatc@ruled
  \def\@fs@pre{\hrule height.8pt depth0pt \kern2pt}%
  \def\@fs@post{\kern2pt\hrule\relax}%
  \def\@fs@mid{\kern2pt\hrule\kern\baselineskip}% original: \kern2pt\hrule\kern2pt
  \let\@fs@iftopcapt\iftrue}
\renewcommand{\@algocf@capt@plain}{above}% formerly {bottom}
\makeatother

\newcommand*{\finalnamedelim}{%
   \ifnumgreater{\value{liststop}}{2}{\finalandcomma}{}%
   \addspace\&\space}

\usepackage[bottom]{footmisc} 

%% Espaciado
\usepackage{setspace}
%\singlespacing        %% 1-spacing (default)
\onehalfspacing       %% 1,5-spacing
%\doublespacing        %% 2-spacing
\setlength\parindent{0pt} 

%%% Norma %%%
%\newcommand{\norm}[1]{\left\lVert#1\right\rVert}

%%% Valor absoluto %%%
\usepackage{mathtools}
\DeclarePairedDelimiter\abs{\lvert}{\rvert}%
\DeclarePairedDelimiter\norm{\lVert}{\rVert}%
% Swap the definition of \abs* and \norm*, so that \abs
% and \norm resizes the size of the brackets, and the 
% starred version does not.
\makeatletter
\let\oldabs\abs
\def\abs{\@ifstar{\oldabs}{\oldabs*}}
%
\let\oldnorm\norm
\def\norm{\@ifstar{\oldnorm}{\oldnorm*}}
\makeatother
%%% %%%

%%%% R code %%%%
\usepackage{listingsutf8}
\usepackage[usenames,dvipsnames]{color}   
\lstset{ 
  language=R,                     % the language of the code
  basicstyle=\footnotesize \ttfamily,  % the size of the fonts that are used for the code
  numbers=left,                   % where to put the line-numbers
  numberstyle=\tiny\color{Blue},  % the style that is used for the line-numbers
  stepnumber=1,                   % the step between two line-numbers. If it is 1, each line
                                  % will be numbered
  numbersep=5pt,                  % how far the line-numbers are from the code
  backgroundcolor=\color{white} ,,  % choose the background color. You must add \usepackage{color}
  showspaces=false,               % show spaces adding particular underscores
  showstringspaces=false,         % underline spaces within strings
  showtabs=false,                 % show tabs within strings adding particular underscores
  frame=single,                   % adds a frame around the code
  rulecolor=\color{black},        % if not set, the frame-color may be changed on line-breaks within not-black text (e.g. commens (green here))
  tabsize=2,                      % sets default tabsize to 2 spaces
  captionpos=b,                   % sets the caption-position to bottom
  breaklines=false,                % sets automatic line breaking
  breakatwhitespace=false,        % sets if automatic breaks should only happen at whitespace
  keywordstyle=\ttfamily,
  commentstyle=\color{RoyalBlue},   % comment style
  stringstyle=\color{ForestGreen},      % string literal style
  inputencoding=utf8/latin1
} 

\newtheorem{defin}{Definici\'on}
\newtheorem*{defin*}{Definici\'on}
\newtheorem*{obs*}{Observaci\'on}
\newtheorem*{theorem*}{Teorema}
\newtheorem*{prop*}{Propiedad}

\newenvironment{dedication}
  {\clearpage           % we want a new page          
   \thispagestyle{empty}% no header and footer
   \vspace*{\stretch{1}}% some space at the top
   \itshape             % the text is in italics
   \raggedleft          % flush to the right margin
  }
  {\par % end the paragraph
   \vspace{\stretch{3}} % space at bottom is three times that at the top
   \clearpage           % finish off the page
  }

%%%%%%%%%%%%%%%%%%%%%%%%%%%%%%%%%%%%%%%%%%%%%%%%%%%%%%%%%%%%%
%%  Documento
%%%%%%%%%%%%%%%%%%%%%%%%%%%%%%%%%%%%%%%%%%%%%%%%%%%%%%%%%%%%%
\begin{document}

\pagestyle{empty} %No headings for the first pages.

\title{UN MODELO BAYESIANO Y NO PARAM\'ETRICO DE REGRESI\'ON SOBRE CUANTILES} %Con este nombre se guardar\'a el proyecto en writeLaTex

%%%%%%%%%%%%%%%%%%%%%%%%%%%%%%%%%%%%%%%%%%%%%%%%%%%%%%%%%%%%%
%%    Secciones iniciales
%%%%%%%%%%%%%%%%%%%%%%%%%%%%%%%%%%%%%%%%%%%%%%%%%%%%%%%%%%%%%

% COVER PAGE
\begin{titlepage}
\begin{center}

\large{INSTITUTO TECNOLOGICO AUTONOMO DE MEXICO}\\[2em]
\hline

%Figura
\begin{center}
	\includegraphics{Figures/Miscellaneous/logo-ITAM.pdf}
\end{center}

\vspace{1em}

{\sc {\bf \large UN MODELO BAYESIANO Y NO PARAMETRICO DE REGRESION SOBRE CUANTILES}}\\[3em]

\textsc{\large TESIS}\\[1em]

\textsc{\normalsize QUE PARA OBTENER EL TITULO DE}\\[1em]

\textsc{\normalsize LICENCIADO EN MATEMATICAS APLICADAS}\\[1em]

\textsc{\normalsize PRESENTA}\\[1em]

\textsc{\Large CARLOS OMAR PARDO GOMEZ}\\[3em]

\end{center}

\vspace*{\fill}
\textsc{CIUDAD DE MEXICO \hspace*{\fill} 2018}

\end{titlepage}
\input{Sections/FirstPage.tex}

\thispagestyle{empty}
\vspace*{\fill}
\begingroup
Con fundamento en los art\'iculos 21 y 27 de la Ley Federal del Derecho de Autor y como titular de los derechos moral y patrimonial de la obra titulada "\textbf{UN MODELO BAYESIANO Y NO PARAM\'ETRICO DE REGRESI\'ON SOBRE CUANTILES}", otorgo de manera gratuita y permanente al Instituto Tecnol\'ogico Aut\'onomo de M\'exico y a la Biblioteca Ra\'ul Baill\'eres Jr., la autorizaci\'on para que fijen la obra en cualquier medio, incluido el electr\'onico, y la divulguen entre sus usuarios, profesores, estudiantes o terceras personas, sin que pueda percibir por tal divulgaci\'on una contraprestaci\'on.

\centering

\hspace{3em}

\textsc{Carlos Omar Pardo G\'omez}

\vspace{4em}

\rule[1em]{20em}{0.5pt} % L\'inea para la fecha

\textsc{Fecha}
 
\vspace{6em}

\rule[1em]{20em}{0.5pt} % L\'inea para la firma

\textsc{Firma}

\endgroup
\vspace*{\fill}
\begin{dedication}        
A Mago.
\end{dedication}
%\clearpage
\setcounter{page}{1}
\chapter*{Agradecimientos}
\markboth{AGRADECIMIENTOS}{AGRADECIMIENTOS} % encabezado 
A Martha, porque gran parte de lo que soy, y lo que sueño ser, se lo debo a tu gu\'ia. A Chevo, por el gran esfuerzo que haces todos los d\'ias por los tuyos. A Jany, por ser mi compañera de aventuras. A Carlos, por tu inigualable ejemplo y sabidur\'ia. A Magally, por tu infinito apoyo. A Janda, Rosy, Don Benja, Betty, Marina, Benja, Renata, Benjam\'in, Blanquita, Luis y Rebeca, por conformar a nuestra pequeña, pero unida familia.

Al Dr. Mart\'inez Ovando, por su confianza y asesor\'ia para sacar adelante esta tesis. Al Dr. Nieto, por introducirme al mundo de la estad\'istica. Al Dr. Mendoza, porque gracias a \'el las palabras \textit{regresi\'on} y \textit{Bayesiano} aparecen en este trabajo. Al Dr. Horacio, por mostrarme c\'omo llevar el conocimiento te\'orico al mercado laboral. Y a los cuatro, por el esfuerzo que le dedicaron a la revisi\'on de este documento.

A Farf\'an, Sol, Alexa, Ari, Rodo, Shaq, Norman, Cisneros, y Amor, porque despu\'es de tantos años de amistad, los considero una extensi\'on de mi familia. A Mago, por tu legado de bondad e inteligencia. A tu familia, por el gran ejemplo de fortaleza. 

A Lucy, por los debrayes matem\'aticos y no matem\'aticos. A Lobato, Wille y Josman, por las m\'ultiples facetas que desarrollamos juntos. A Alfredo, Joaquín e Isaac, porque estar rodeado de gente tan inteligente me hizo ser mejor. A Luis y Lalo, por llevar nuestra amistad m\'as all\'a de la escuela. Y a todos con los que compart\'i clases en la carrera de Matem\'aticas Aplicadas.

A Pau Chambon, por aceptarme en todas mis facetas. A Jaime y Rubén, por las risas. A La Bambi, porque me ayudaron a derribar muchos de mis prejuicios. A Fer Langarica y Fer Zepeda, porque estudiar Actuar\'ia me dio excelentes amistades.

A Alain, Cass, y Rafa, por la locura que convertimos en realidad. A Belén, Mariana, Adri, Andrea, Naty y Saus, por permitirme guiarlos. A Brenda, por tu intensidad. A Pato y a Fer Landa, por compartir tan buenas charlas. Y, en general, a todos los miembros de Doxia y SALITAM 2015. 

A Omar Trejo, Luis Garc\'ia, Jorge Alag\'on, Abraham Nieto y Conchita Tapia, por confiar en mi capacidad laboral. A Mois\'es Arizpe, a Mario Guerrero, y, en general, a los equipos de Datata, Kantar, BBVA Bancomer y Datio, por lo que han aportado a mi desarrollo profesional.

A Aar\'on, Fer Rubio y Fer Hern\'andez, por nuestras confidencias. A Vann, por lo que compartimos. A Jess, por siempre sacar mi mejor versi\'on. A Irene y Erick Bustamante, por su confianza. A Luis y Napo, por nuestras interesantes conversaciones. A Mary, por tu risa y tu sonrisa. Y a todos los que han aportado a mi crecimiento personal despu\'es de la universidad. 

Finalmente, al ITAM, por una eduaci\'on de tan alta calidad. A quienes intervinieron para que pudiera disfrutar de las becas Baillères y Mancera. A mis exalumnos, porque disfrut\'e mucho compartir conocimiento. Y a los profesores que aportaron a mi crecimiento, entre los que destaco a Guillermo Grabinsky, Zeferino Parada, Felipe Gonz\'alez, Alberto Tubilla, Marcela Gonz\'alez, Sergio Garc\'ia, Isaac Katz, Valeria Moy y Rosario Sarmiento.
\clearpage
\setcounter{page}{1}
\chapter*{Agradecimientos}
\markboth{AGRADECIMIENTOS}{AGRADECIMIENTOS} % encabezado 
¡Muchas gracias a todos!
\pagestyle{plain}
\chapter*{Prefacio}
\markboth{PREFACIO}{PREFACIO} % encabezado 

%%  Contenido
\addtocontents{toc}{\protect\vspace*{\baselineskip}}

El tema de esta tesis es describir un modelo de regresi\'on sobre cuantiles, debido a diversas bondades que presenta sobre el tradicional an\'alisis de regresi\'on a la media. Adem\'as, el modelo tiene una construcci\'on probabil\'istica realizada a través del paradigma Bayesiano y presenta una gran flexibilidad, al contar con componentes no param\'etricos. Asimismo, en este trabajo se abordan los modelos tradicionales de regresi\'on, para entender las desventajas contrarrestadas por el nuevo modelo.

El cap\'itulo 1 describe la importancia de las aproximaciones distintas al modelo tradicional de regresi\'on (a la media, lineal y con error normal), así como la evoluci\'on hist\'orica de estos modelos alternativos. El cap\'itulo 2 introduce al paradigma Bayesiano y sus fundamentos generales. El cap\'itulo 3 se centra en los modelos Bayesianos tradicionales de regresi\'on, tanto a la media, como sobre cuantiles. El cap\'itulo 4 plantea la especificaciones no param\'etricas para la tendencia y el error. El cap\'itulo 5 define el modelo propuesto por esta tesis y explica el algoritmo necesario para realizar inferencia y predicci\'on. El cap\'itulo 6 muestra simulaciones y aplicaciones del modelo. Finalmente, el cap\'itulo 7 hace referencia a las conclusiones de esta tesis, adem\'as de describir el trabajo futuro que se podr\'ia desarrollar.

%%%%%%%%%%%%%%%%%%%%%%%%%%%%%%%%%%%%%%%%%%%%%%%%%%%%%%%%%%%%%
%%    Índice
%%%%%%%%%%%%%%%%%%%%%%%%%%%%%%%%%%%%%%%%%%%%%%%%%%%%%%%%%%%%%

\pagestyle{empty}
\begingroup
\setlength{\parskip}{0pt}
\tableofcontents
\endgroup

%%%%%%%%%%%%%%%%%%%%%%%%%%%%%%%%%%%%%%%%%%%%%%%%%%%%%%%%%%%%%
%%    Capitulos
%%%%%%%%%%%%%%%%%%%%%%%%%%%%%%%%%%%%%%%%%%%%%%%%%%%%%%%%%%%%%

\clearpage
\pagestyle{fancy}
\fancyhf{}
\fancyhead[R]{\nouppercase{\leftmark}}
\cfoot{\thepage}

\chapter{Introducci\'on}

Detr\'as de cualquier modelo de regresi\'on la intenci\'on es entender alguna caracter\'istica asociada con una variable aleatoria, en funci\'on de un conjunto de variables potencialmente explicativas o predictivas para tal caracter\'istica. Ha sido com\'un resumir esta dependencia mediante alguna medida de tendencia central, condicionada a los valores de las covariables.

La medida de tendencia central tradicionalmente usada ha sido la media, dando lugar a los modelos de \textit{regresi\'on a la media}, con sus variantes lineal y no lineal, simple y m\'ultiple, con error normal y no normal. Este tipo de modelos tiene un buen n\'umero de ventajas, entre las que destacan el bajo costo de estimaci\'on y la facilidad de interpretaci\'on de sus par\'ametros (principalmente en la versi\'on lineal). Sin embargo, como mencionan \cite{Hao_FrecQuantReg}, tienen tres grandes limitaciones, que a continuaci\'on se mencionan.

La primera es que, generalmente, la estimaci\'on del modelo busca minimizar la diferencia entre el valor esperado te\'orico y la media observada. Por lo tanto, la inferencia sobre los valores lejanos a la media, que suelen ser de inter\'es en ciertos contextos, como los seguros o las finanzas, puede ser inexacta.

La segunda es que algunos fen\'omenos de estudio tienen distribuciones de colas pesadas, principalmente en las ciencias sociales. Esto da lugar a valores at\'ipicos, mismos que pueden sesgar estimaciones de la media, mientras prácticamente no afectan a otros estadísticos, como la mediana.

La tercera es que al focalizar la intervenci\'on de las variables explicativas \'unicamente en la media, y darle una distribución predefinida al error, se suelen dejar de lado caracter\'isticas propias del fen\'omeno que se est\'a estudiando. Por ejemplo, cada cuantil, adem\'as de diferir en el nivel, podr\'ia tener una forma funcional distinta.

Debido a esto, desde mitades del siglo XVIII han surgido alternativas a este tipo de modelos. Seg\'un \cite{Hao_FrecQuantReg}, la primera conocida data de 1760, cuando el jesuita croata Rudjer Josip Boscovich visit\'o Londres en b\'usqueda de consejo computacional para su novedoso modelo de \textit{regresi\’on a la mediana}. De nueva cuenta se busc\'o una medida de tendencia central, pero con otras bondades. Por ejemplo, ser una mejor medida informativa para distribuciones asim\'etricas y menos susceptible a valores at\'ipicos. 

As\'i como los modelos de regresi\'on a la media son com\'unmente relacionados con la minimizaci\'on de los errores cuadr\'aticos, los modelos de regresi\'on a la mediana lo son con la minimizaci\'on de los errores absolutos. Debido a la no diferenciabilidad, tuvieron que pasar muchos años para que lograran ser viables, hasta que el poder computacional y los algoritmos de programaci\'on lineal lo permitieron.

Cabe recordar que, a grandes rasgos y sin dar a\'un una definici\'on formal, el cuantil $p$-\'esimo es aquel valor tal que el $p \times 100\%$ de los valores est\'an por debajo de \'el, y el $(1-p)\times 100\%$, por encima. As\'i, la mediana es un caso particular de un cuantil, espec\'ificamente el $0.5$-\'esimo. Esto abre la idea de que otros cuantiles tambi\'en podr\'ian ser modelados en funci\'on de variables explicativas y no necesariamente tienen que ser una medida de tendencia central. 

Los \textit{modelos de regresi\'on sobre cuantiles} fueron introducidos por \cite{Koenker_QuantReg}, y han permitido concentrarse en la estimaci\'on de valores de inter\'es para los modeladores, sin importar que est\'en alejados de la media. Adem\'as, el c\'alculo de diversos cuantiles para un mismo fen\'omeno ha permitido entender mejor la forma y propiedades de las distribuciones condicionales de la variable de respuesta.

El paradigma Bayesiano ha desarrollado este tipo de modelos de forma posterior a otros paradigmas. \cite{Walker_BayesAccFail}, \cite{Kottas_BaySemiparamMed} y \cite{Hanson_PolyaTrees} desarrollaron modelos para la mediana, suponiendo una distribuci\'on no param\'etrica del error. \cite{Yu_BayQuantReg} y \cite{Tsionas_BayQuantInf} desarrollaron inferencia param\'etrica, basados en la distribuci\'on asim\'etrica de Laplace para los errores. Por otro lado, \cite{Lavine_LikeQuant} y \cite{Dunson_ApproxBayes} usaron una perspectiva distinta y propusieron una aproximaci\'on de la verosimilitud para cuantiles.

Por la misma naturaleza compleja del nuevo desaf\'io, los modelos de regresi\'on sobre cuantiles han retomado un concepto que tambi\'en ha tenido auge en los de regresi\'on a la media: la distribuci\'on no param\'etrica de los errores, misma que generaliza la idea tradicional del error normal. 

Desafortunadamente, a diferencia de los modelos de regresi\'on a la media, que ya han propuesto alternativas para este tema, los modelos Bayesianos de regresi\'on sobre cuantiles han hecho muy poco por romper con la relaci\'on lineal en los par\'ametros entre la variable de respuesta y las covariables. Y casi todos aquellos que lo han logrado, han tenido que recurrir a estimaciones no probabil\'isticas o no Bayesianas, para resolver alguna parte del problema.

Esta tesis tiene la finalidad de sustituir simult\'aneamente las ideas tradicionales de regresi\'on a la media, linearidad y distribuci\'on normal de los errores, por un enfoque m\'as flexible, y que se mantenga totalmente probabil\'istico. Para ello, rescata las ideas de \cite{Kottas_NotParamQuantReg} y \cite{Kottas_SemiparamQuantReg}, proponiendo un modelo Bayesiano y no param\'etrico, \'util en el contexto de regresi\'on sobre cuantiles.

\newpage
\chapter[Paradigma Bayesiano]{Paradigma bayesiano\footnote{Las ideas de este cap\'itulo son retomadas de \cite{Denison_BayesMethods}.}\textsuperscript{,}\footnote{Esta tesis da como aceptados los axiomas de coherencia de la Teoría de la Decisión, mismos que pueden ser encontrados, por ejemplo, en \cite{Fishburn_Axioms}. Por lo tanto, entiende al paradigma Bayesiano como el coherente para hacer estad\'istica, cuando una toma de decisi\'on con incertidumbre es el objetivo final del estudio. 
}}
\label{chap:Bayesian}

\section{Inferencia de variables aleatorias}

Un problema clásico de la estad\'istica es el de hacer predicci\'on, utilizando la informaci\'on de los datos que ya han sido observados. Por ejemplo, es posible pensar que ya se tiene el conjunto de $n$ datos observados $\{y_1, \ldots, y_n\}$ y se desea hacer predicci\'on acerca del valor del dato $y_{n+1}$, que a\'un no ha sido observado. Para esto, se podr\'ia usar la probabilidad condicional
\begin{equation*}
    \mathbb{P}(y_{n+1}|y_1,\ldots,y_n) =
    \frac{\mathbb{P}(y_{n+1} \cap \{y_1, \ldots, y_n\})}{\mathbb{P}(y_1, \ldots, y_n)} =
    \frac{\mathbb{P}(y_1, \ldots, y_n,y_{n+1})}{\mathbb{P}(y_1, \ldots, y_n)},
\end{equation*}
pero esto requerir\'ia conocer la funci\'on conjunta, misma que puede ser compleja por la estructura de dependencia de los datos.\footnote{En este trabajo se usar\'a la notaci\'on $\mathbb{P}$ como una forma general de definir una medida de probabilidad, independientemente de los detalles te\'oricos sobre an\'alisis y medibilidad.}

Este problema puede ser abordado mediante el uso del Teorema de representaci\'on general de de Finetti. Para ello, antes se dar\'a una definici\'on.

\begin{defin*}
    Sea $(y_1,y_2,\ldots)$, una sucesi\'on de variables aleatorias, cuya distribuci\'on de probabilidad conjunta est\'a dada por $\mathbb{P}(y_1,y_2,\ldots)$. Sea $\psi$ una funci\'on biyectiva que crea una permutaci\'on finita del conjunto $\{1,2,\ldots\}$, es decir, permuta un n\'umero finito de elementos y al resto los deja fijos.  
    Se dice entonces que $(y_1,y_2,\ldots)$ es una \textbf{sucesi\'on aleatoria infinitamente intercambiable} si se cumple que 
    \begin{equation*}
        \mathbb{P}(y_1,y_2,\ldots) = \mathbb{P}(y_{\psi(1)},y_{\psi(2)},\ldots),
    \end{equation*}
    para cualquier permutaci\'on $\psi$.
\end{defin*}

En pocas palabras, una sucesi\'on $(y_1,y_2,\ldots)$ se considerar\'a infinitamente intercambiable si el orden en que se etiquetan las variables no afecta su distribuci\'on conjunta. Es importante hacer notar que la com\'unmente usada independencia implica intercambiabilidad, pero lo contrario no se cumple. Es decir, la intercambiabilidad es un supuesto menos r\'igido que la independencia.

Dicho esto, es momento de plantear el \textbf{\textit{Teorema de represesentaci\'on general de de Finetti}}.\footnote{Una demostraci\'on de este teorema puede ser encontrada en \cite{Schervish_TheoryStats}.}

\begin{theorem*}
    Sea $(y_1,y_2,\ldots)$ una sucesi\'on aleatoria infinitamente intercambiable de valores reales. Entonces existe una distribuci\'on de probabilidad $F$ sobre $\mathcal{F}$, el espacio de todas las distribuciones, de forma que la probabilidad conjunta de $(y_1,y_2,\ldots)$ se puede expresar como
    \begin{equation*}
        \mathbb{P}(y_1,y_2,\ldots) =
        \int_{\mathcal{F}}\left[\prod_{k=1}^\infty \mathbb{P}(y_k|G)\right]dF(G),
    \end{equation*}
    con
    \begin{equation*}
        F(G) = \lim_{n \to \infty} F(G_n),
    \end{equation*}
    donde $F(G_n)$ es una funci\'on de distribuci\'on evaluada en la funci\'on de distribuci\'on emp\'irica definida por
    \begin{equation*}
        G_n(y) = \frac{1}{n} \sum_{i=1}^n I(y_i \leq y).
    \end{equation*}
    En otras palabras, el Teorema de de Finetti dice que $\{y_1,y_2,\ldots\}$ es un conjunto de variables aleatorias condicionalmente independientes, dado que se supone provienen de la distribuci\'on $G$. Al mismo tiempo, $G$ es una realizaci\'on de otra variable aleatoria cuya distribuci\'on est\'a dada por $F(G)$.
\end{theorem*}

Cabe hacer notar que dicho teorema plantea la distribuci\'on conjunta de $(y_1,y_2,\ldots)$ como una mezcla de verosimilitudes condicionalmente independientes en $G$, donde el peso asociada a cada una depende de $F(G)$. Por lo tanto, $F(G)$ expresa la incertidumbre acerca de cu\'an probable es que $G$ sea id\'onea para explicar el fen\'omeno, a\'un sin observar los datos.

Un subconjunto del espacio de todas las distribuciones $\mathcal{F}$ es el espacio de las distribuciones param\'etricas, es decir, aquellas que pueden ser descritas en su totalidad \'unicamente señalando el valor de un vector de par\'ametros $\theta$ de tamaño finito, mismo que puede tomar valores en todo un soporte $\Theta$\footnote{En lo que resta de esta secci\'on las proposiciones y resultados har\'an referencia al conjunto de distribuciones param\'etricas, apelando a que as\'i los nuevos conceptos tengan mayor claridad para el lector, considerando que este tipo de distribuciones son con las que se suele estar m\'as familiarizado. Todos ser\'an generalizables al conjunto de distribuciones $\mathcal{F}$.}. Por lo tanto, si se hace el supuesto adicional que la distribuci\'on marginal de $y_i$ es param\'etrica, con vector de par\'ametros desconocido $\theta$, se obtiene como corolario del Teorema de de Finetti que
\begin{equation*}
    \mathbb{P}(y_1,y_2,\ldots) =
    \int_{\Theta}\left[\prod_{k=1}^\infty \mathbb{P}(y_k|\theta)\right]\mathbb{P}(\theta)d\theta.
\end{equation*}

Siguiendo el razonamiento anterior, $\mathbb{P}(\theta)$ refleja la incertidumbre del modelador acerca de cu\'al es el vector $\theta$ que origin\'o los datos, antes de observarlos. 

La intuici\'on detr\'as de este resultado es que, al igual que en otros paradigmas, se supone a un s\'olo vector $\theta$ como aquel del que provinieron los datos, pero es desconocido, y la tarea es estimarlo. Una particularidad del paradigma Bayesiano es expresar la incertidumbre que tiene el modelador acerca del valor verdadero mediante la asignaci\'on de una distribuci\'on de probabilidad para $\theta$, sujeta a la informaci\'on inicial o conocimiento previo a observar los datos $(CP)$ que se tenga del fen\'omeno. Es decir, $\mathbb{P}(\theta|CP)$. Como una simplificaci\'on de la notaci\'on, en la literatura normalmente se escribe $\mathbb{P}(\theta) = \mathbb{P}(\theta|CP)$ y se conoce como la \textit{probabilidad inicial} del par\'ametro.

As\'i, el Teorema de representaci\'on general garantiza que para variables aleatorias infinitamente intercambiables existe una distribuci\'on del vector de par\'ametros, tal que la probabilidad conjunta se puede expresar como la verosimilitud de variables condicionalmente independientes, dado un vector de par\'ametros, multiplicada por la probabilidad de que dicho vector de par\'ametros sea del que efectivamente provinieron las observaciones. Con el teorema se prueba la existencia de dicha representaci\'on, mas no su unicidad. En los siguientes p\'arrafos se har\'a uso de ella para realizar inferencia en variables aleatorias.

Regresando al problema inicial, y bajo los supuestos reci\'en mencionados, es posible escribir
\begin{equation*}
    \mathbb{P}(y_{n+1} | y_1,\ldots,y_n) =
    \int_{\Theta} 
    \mathbb{P}(y_{n+1}|\theta) \mathbb{P}(\theta|y_1,\ldots,y_n)d\theta,
\end{equation*}
donde a su vez, usando el \textbf{Teorema de Bayes}, se obtiene que
\begin{equation*}
\begin{aligned}
    \mathbb{P}(\theta|y_1,\ldots,y_n) &=
    \frac{\mathbb{P}(y_1,\ldots,y_n|\theta)\times\mathbb{P}(\theta)}
    {\mathbb{P}(y_1,\ldots,y_n)},
\end{aligned}
\end{equation*}
que en el paradigma Bayesiano se conoce como la \textit{probabilidad posterior} del par\'ametro. 

Se puede observar que el denominador es constante respecto a $\theta$, por lo que usualmente la probabilidad condicional del vector de par\'ametros no se expresa como una igualdad, sino con la expresi\'on
\begin{equation*}
    \mathbb{P}(\theta|y_1,\ldots,y_n) 
    \propto 
    \mathbb{P}(y_1,\ldots,y_n|\theta) \times \mathbb{P}(\theta),
\end{equation*}
y s\'olo difiere de la igualdad por una constante que permita que, al integrar sobre todo el soporte de $\theta$, el resultado sea igual a 1.

Cabe resaltar que debido a la independencia condicional dada por el Teorema de de Finetti, el factor de verosimilitud puede ser reescrito como
\begin{equation*}
    \mathbb{P}(y_1,\ldots,y_n|\theta)  = \prod_{i=1}^n \mathbb{P}(y_i|\theta).
\end{equation*}
Por lo tanto, se confirma que el aprendizaje en el paradigma Bayesiano se obtiene como
\begin{equation*}
    Posterior \propto Verosimilitud \times Inicial.
\end{equation*}
Es decir, el conocimiento final surge de conjuntar el conocimiento inicial con la informaci\'on contenida en los datos.

Es importante notar que bajo este enfoque se obtiene una distribuci\'on de probabilidad para un valor futuro de $y_{n+1}$. \'Esta se puede utilizar para el c\'alculo de predicciones puntuales o intervalos (que en el caso del paradigma Bayesiano son de \textit{probabilidad}), mediante el uso de la Teor\'ia de la Decisi\'on y funciones de utilidad o p\'erdida.

\section{Propiedad conjugada}

En los casos en los que la distribuci\'on posterior de los par\'ametros tiene la misma forma funcional que la distribuci\'on inicial, se dice que la distribuci\'on de los par\'ametros pertenece a una \textbf{familia param\'etrica conjugada}, para la verosimilitud respectiva.

Esta propiedad es conveniente, porque permite a la distribuci\'on posterior tener forma anal\'itica cerrada, evitando tener que usar m\'etodos num\'ericos para aproximarla. Adem\'as permite ver de forma m\'as clara c\'omo afectan los datos a la actualizaci\'on, respecto a la distribuci\'on inicial.

Sin embargo, el rango de posibles modelos conjugados puede resultar limitado en algunos contextos pr\'acticos debido a que el fen\'omeno en estudio puede ser mejor representado con otras distribuciones, que usualmente no pertenecen a familias conjugadas.

\section[Inferencia con variables explicativas]{Inferencia con variables explicativas\raisebox{.3\baselineskip}{\normalsize\footnotemark}}\footnotetext{Esta secci\'on carece de una formalidad completa, pero busca darle la intuici\'on al lector para generalizar el resultado del Teorema de de Finetti en el contexto en el que se desarrollar\'a este trabajo. Para una explicaci\'on m\'as formal, consultar \cite{Dawid_Exchangeability}.}

Como se ver\'a en el siguiente cap\'itulo de esta tesis, un problema com\'un es estimar la distribuci\'on de cierta sucesi\'on de variables aleatorias $(y_1,y_2,\ldots)$, condicionada a las observaciones $(x_1,x_2,\ldots)$ de otras variables usualmente llamadas explicativas o predictivas (cada $x_i$ es un vector de dimensi\'on finita). En este caso, la sucesi\'on $(y_1,y_2,\ldots)$ ya no es intercambiable, porque cada valor $y_i$ depende del valor de su respectiva $x_i$, por lo que no es posible aplicar de manera directa el Teorema de de Finetti. Para hacerlo de manera indirecta, se introducir\'a el concepto de intercambiabilidad parcial.

\begin{defin*}
    Sea $(y_1,y_2,...)$ dada $(x_1,x_2,...)$, una sucesi\'on numerable de variables aleatorias, asociada con los correspondientes valores de sus variables explicativas, y cuya distribuci\'on de probabilidad conjunta est\'a dada por $\mathbb{P}(y_1,y_2,\ldots|x_1,x_2,\ldots)$. Sea $\psi$ una funci\'on biyectiva que crea una permutaci\'on finita de un conjunto, es decir, permuta un n\'umero finito de elementos y al resto los deja fijos. Sean $\tilde{x}_1,\tilde{x}_2,\ldots$ los distintos valores \'unicos que toman las $x$'s, y $y_{k_i}$ una $y$, tal que el valor de su correspondiente variable explicativa es $\tilde{x}_k$.  
    Entonces, se dice que $(y_1,y_2,\ldots)$ es una \textbf{sucesi\'on aleatoria infinita y parcialmente intercambiable} si se cumple que 
    \begin{equation*}
        \mathbb{P}(y_{k_1},y_{k_2},\ldots|\tilde{x}_k) = \mathbb{P}(y_{\psi(k_1)},y_{\psi(k_2)},\ldots|\tilde{x}_k),
    \end{equation*}
    para cualquier permutaci\'on $\psi$ y para todos los diferentes valores $k$.
    
    Es decir, todas aquellas $y$'s cuyas $x$'s tienen el mismo valor, son infinitamente intercambiables entre s\'i. 
\end{defin*}

Si adem\'as se cumpliera que el orden de los valores \'unicos $\tilde{x}$'s es intercambiable en el sentido de de Finetti, entonces, intuitivamente se podr\'ia tomar la $G$ del Teorema de de Finetti como dependiente de las $x_i$'s, $G(x_i)$, y se obtendr\'ia que
\begin{equation*}
    \mathbb{P}(y_1,y_2,\ldots|x_1,x_2,\ldots) =
    \int_{\mathcal{F}}\left[\prod_{i=1}^\infty \mathbb{P}(y_i|G(x_i))\right]dF(G(x_1,x_2,\ldots)),
\end{equation*}
donde las $y$'s resultar\'ian ser independientes entre s\'i, condicionadas a una distribuci\'on que depende de la $x_i$ asociada.

El tema del siguiente cap\'itulo ser\'a la discusi\'on de m\'etodos de inferencia sobre las variables $y$, dentro de este contexto espec\'ifico de dependencia de una variable explicativa o predictiva $x$, normalmente conocidos como \textbf{modelos de regresi\'on}.

\newpage
\chapter[Modelos de regresi\'on]{Modelos de regresi\'on}

\section{Concepto general}

Los modelos de regresi\'on tienen como objetivo describir la distribuci\'on de una variable aleatoria $y \in \mathbb{R}$, generalmente llamada \textit{variable de respuesta}, condicional a los valores de las variables $x \in \mathbb{R}^n$, conocidas como \textit{covariables} o \textit{variables de entrada}. Visto en t\'erminos matem\'aticos, se puede expresar como
\begin{equation*}
    y|x \sim \mathbb{P}(y|x).
\end{equation*}

Si bien esta relaci\'on se da por hecha y es fija, normalmente es desconocida. Por lo tanto, la intenci\'on de estos modelos es realizar alguna aproximaci\'on de ella. Dado que es complicado aproximar con exactitud toda la distribuci\'on, com\'unmente se utilizan un n\'umero finito de par\'ametros para describirla. Adem\'as, la interpretaci\'on de dichos par\'ametros suele tener relevancia para el modelador, como es el caso de la media o la mediana. 

Tambi\'en es importante resaltar que este trabajo tiene inter\'es en modelar a $y$, dado que ya se observaron los valores de $x$. Sin embargo, se podr\'ia pensar en modelos donde haga sentido la distribuci\'on conjunta de $y$ y $x$, misma que se podr\'ia obtener como
\begin{equation*}
    \mathbb{P}(y,x) = \mathbb{P}(y|x) \times \mathbb{P}(x).
\end{equation*}

\section{Regresión a la media}

La \textit{regresi\'on a la media} es el caso particular m\'as usado de los modelos de regresi\'on, tanto en el paradigma Bayesiano, como en otros. Esto sucede debido al bajo uso de recursos que requiere su estimaci\'on, particularmente porque familias conjugadas suelen hacer sentido en este contexto. De hecho, fuera del paradigma Bayesiano suelen ser com\'unmente usados debido a que se asocian con la minimizaci\'on de t\'erminos cuadrados, con las bondades de diferenciabilidad que eso implica. Adem\'as se valora la capacidad interpretativa que suelen tener los par\'ametros estimados.

En notaci\'on probabil\'istica, retomando el hecho de que $y|x \sim \mathbb{P}(y|x)$, el modelo de regresi\'on a la media busca aproximar a la funci\'on $f$, tal que 
\begin{equation*}
    \mathbb{E}(y|x) = f(x).
\end{equation*}
Para hacer esto, normalmente se vale del supuesto que
\begin{equation*}
    y = f(x) + \varepsilon,
\end{equation*}
con $\varepsilon \in \mathbb{R} \sim \mathcal{N}(0,\sigma^2)$ (denominado com\'unmente como el \textit{error aleatorio}), y siendo $f: \mathbb{R}^n \rightarrow \mathbb{R}$ y $\sigma^2 \in \mathbb{R_+}$ desconocidas, de forma que
\begin{equation*}
    y | x, f, \sigma^2 \sim \mathcal{N}(f(x),\sigma^2).
\end{equation*}

Adem\'as, se supone independencia entre $\varepsilon$s, es decir, para todo $\tilde{\varepsilon} \neq \dot{\varepsilon}$, $\tilde{\varepsilon}$ y $\dot{\varepsilon}$ son independientes. Por lo tanto, sean $\tilde{x}$ las covariables asociadas a la variable de respuesta $\tilde{y}$, y $\dot{x}$ las asociadas a $\dot{y}$, se tiene que $\tilde{y} | \tilde{x}, f, \sigma^2$ es condicionalmente independiente a $\dot{y} | \dot{x}, f, \sigma^2$.

\subsection[Modelo tradicional]{
    Modelo tradicional
    \footnote{Algunas ideas de esta subsecci\'on son retomadas de \cite{Denison_BayesMethods} y \cite{Bannerjee_BayLinMod}.}
}

La \textit{regresi\'on lineal a la media} es el caso particular m\'as usado en el contexto de \textit{regresi\'on a la media}. Consiste en definir
\begin{equation*}
    f(x) = x^T\beta,
\end{equation*}
donde $\beta \in \mathbb{R}^n$ se piensa con valores constantes, pero desconocidos, y la tarea es estimarlos, al igual que $\sigma^2$.

Para hacer esto, el enfoque Bayesiano le asigna una distribución inicial de probabilidad a ambos par\'ametros, reflejando la incertidumbre que tiene el modelador acerca de su valor real. Es decir, sea $CP$ la hip\'otesis o el conocimiento previo al que tiene acceso el modelador, se tiene que 
\begin{equation*}
    \beta,\sigma^2 \sim \mathbb{P}(\beta,\sigma^2|CP).
\end{equation*}

Como ya se mencion\'o en el cap\'itulo anterior, se omitir\'a escribir la distribuci\'on condicional respecto a $CP$ por simplificaci\'on de la notaci\'on, pero es importante no olvidar su existencia.

Sea $\{(x_i,y_i)| x_i \in \mathbb{R}^n, y_i \in \mathbb{R}, i \in \{1,...,m\} \}$ el conjunto de datos observados de las variables de respuesta y de las covariables. Es posible representar este mismo conjunto con la notaci\'on matricial $\{X,Y | X \in \mathbb{R}^{m \times n}, Y \in \mathbb{R}^m\}$. \footnote{Por simplificaci\'on y limpieza de notaci\'on en este trabajo se escribir\'an de igual manera variables aleatorias y los datos en efecto observados, considerando que en cada caso el contexto ser\'a suficiente para saber de cu\'al se est\'a hablando, siendo asociadas las letras min\'usculas a una \'unica observaci\'on y las may\'usuclas a una matriz de observaciones.} Sea $\mathcal{E} \in \mathbb{R}^m$ el vector de errores aleatorios, tal que $\mathcal{E} \sim \mathcal{N}(0,\sigma^2 I)$. El modelo se puede reescribir como:
\begin{equation*}
    Y = X\beta + \mathcal{E} \sim \mathcal{N}(X\beta,\sigma^2 I).
\end{equation*}

Por el Teorema de Bayes,
\begin{equation*}
\begin{aligned}
    \mathbb{P}(\beta,\sigma^2 | Y, X) 
    &= \frac{\mathbb{P}(Y| X, \beta, \sigma^2) \times \mathbb{P}(\beta, \sigma^2 | X)}{P(Y | X)} \\
    &= \frac{\mathbb{P}(Y| X, \beta, \sigma^2) \times \mathbb{P}(\beta, \sigma^2)}{\mathbb{P}(Y | X)} \\
    &\propto \mathbb{P}(Y| X, \beta, \sigma^2) \times \mathbb{P}(\beta, \sigma^2), \\
\end{aligned}
\end{equation*}
donde $\mathbb{P}(Y| X, \beta, \sigma^2)$ es la verosimilitud de los datos observados y, debido a la independencia condicional, se puede calcular como $\mathbb{P}(Y| X, \beta, \sigma^2) = \mathcal{N}(X\beta,\sigma^2 I) = \prod_{i=1}^m \mathcal{N}(x_i^T\beta,\sigma^2)$. Por otro lado, $\mathbb{P}(\beta,\sigma^2)$ es la distribuci\'on inicial de los par\'ametros.

Por conveniencia anal\'itica, hay una distribuci\'on inicial com\'unmente usada para los par\'ametros $\beta$ y $\sigma$ debido a que es conjugada respecto a la distribuci\'on Normal de los datos. Su nombre es \textit{Normal-Gamma Inversa (NGI)} y se dice que $\beta,\sigma^2 \sim \mathcal{NGI}(M,V,a,b)$, si
\begin{equation*}
\begin{aligned}
    \mathbb{P}(\beta,\sigma^2) 
    &= \mathbb{P}(\beta|\sigma^2) \times \mathbb{P}(\sigma^2) \\
    &= \mathcal{N}(\beta|M, \sigma^2 V) \times \mathcal{GI}(\sigma ^2|a,b) \\
    &\propto (\sigma^2)^{-(a+(n/2)+1)} \exp\left(-\frac{(\beta-M)^TV^{-1}(\beta-M) + 2b}{2\sigma^2}\right),
\end{aligned}
\end{equation*}
donde $M$ es la media inicial de los coeficientes, $\sigma^2 V$ su varianza, y $a$, $b$ son los par\'ametros iniciales de forma y escala, respectivamente, de $\sigma ^2$. 

Aprovechando la propiedad conjugada, es posible escribir la probabilidad posterior de los par\'ametros como
\begin{equation*}
\begin{aligned}
    \mathbb{P}(\beta,\sigma^2 | Y, X) 
    &\propto \mathbb{P}(Y| X, \beta, \sigma^2) \times \mathbb{P}(\beta, \sigma^2), \\
    &\propto (\sigma^2)^{-(\bar{a}+(n/2)+1)} \exp\left(-\frac{(\beta-\bar{M})^T\bar{V}^{-1}(\beta-\bar{M}) + 2\bar{b}}{2\sigma^2}\right),
\end{aligned}
\end{equation*}
donde
\begin{equation*}
\begin{aligned}
    \bar{M} &= (V^{-1} + X^TX)^{-1} (V^{-1}M + X^TY), \\
    \bar{V} &= (V^{-1} + X^TX)^{-1}, \\
    \bar{a} &= a + n/2, \\
    \bar{b} &= b + \frac{\bar{M}^TV^{-1}M + Y^TY - \bar{M}^T\bar{V}^{-1}\bar{M}}{2}.
\end{aligned}
\end{equation*}

Es decir, la distribuci\'on posterior de $(\beta,\sigma^2)$ es \textit{Normal - Gamma Inversa}, con par\'ametros $\mathcal{NGI}(\bar{M},\bar{V},\bar{a},\bar{b})$.

Si se tiene una nueva matriz de covariables $X_*$ y se desea hacer predicci\'on de las respectivas variables de salida $Y_*$, es posible hacer inferencia con los datos observados como se detalla a continuaci\'on.
\begin{equation*}
\begin{aligned}
    \mathbb{P}(Y_*|X_*,Y,X)
    &= \int \int \mathbb{P}(Y_*|X_*,\beta,\sigma^2) \times \mathbb{P}(\beta,\sigma^2|Y,X) d\sigma^2 d\beta \\
    &= \int \int \mathcal{N}(Y_*|X_*\beta,\sigma^2I) \times \mathbb{P}(\beta,\sigma^2|Y,X) d\sigma^2 d\beta.
\end{aligned}
\end{equation*}

Particularmente, si se contin\'ua con el modelo conjugado \textit{Normal - Gamma Inversa / Normal}, es posible encontrar la soluci\'on anal\'itica:
\begin{equation*}
\begin{aligned}
    \mathbb{P}(Y_*|X_*,Y,X)
    &= \int \int \mathcal{N}(Y_*|X_*\beta,\sigma^2I) \times \mathcal{NGI}(\beta,\sigma^2|\bar{M},\bar{V},\bar{a},\bar{b}) d\sigma^2 d\beta \\
    &= MVSt_{2\bar{a}} 
       \left(
        X_*\bar{M},\frac{\bar{b}}{\bar{a}}\left(I + X_*\bar{V}X_*^T\right)
       \right),
\end{aligned}
\end{equation*}
donde $MVSt$ es la distribuci\'on \textit{t-Student} multivariada, y cuya definici\'on se describe en el \autoref{chap:Distributions}.

\section{Regresión sobre cuantiles}

La \textit{regresi\'on sobre cuantiles} es una alternativa que se ha desarrollado reci\'entemente y que permite enfocarse en aspectos alternativos de la distribuci\'on, como lo que pasa en las colas. Adem\'as relaja supuestos de la regresi\'on a la media, como la simetr\'ia inducida por el error normal.

\begin{defin}
Sea $F_y$ la funci\'on de distribuci\'on de la variable aleatoria $y$, entonces el cuantil $p$-\'esimo de dicha variable aleatoria es aquel valor $q_p$ tal que
\begin{equation*}
    F_y(q_p) = p.
\end{equation*}
Equivalentemente, la funci\'on que regresa el cuantil p-\'esimo de la variable aleatoria $y$ se escribe
\begin{equation*}
    q_p(y) = F_y^{-1}(p),
\end{equation*}
cuando $F_y^{-1}$ est\'a bien definida.
\end{defin}
Dicho en otras palabras, si se tiene un conjunto grande de realizaciones de una variable aleatoria $y$, se esperar\'a que el $p \times 100\%$ est\'e por debajo de $q_p(y)$ y el $(1-p) \times 100\%$ est\'e por arriba. Por ejemplo, la mediana es un caso particular de un cuantil, espec\'ificamente el $0.5$-\'esimo. 

En notaci\'on probabil\'istica, se buscar\'a aproximar a la funci\'on $f$, tal que 
\begin{equation*}
    q_p(y|x) = f(x),
\end{equation*}
para $p \in (0,1)$ fijo arbitrario.

Para hacer esto, normalmente se vale del supuesto que
\begin{equation*}
    y = f_p(x) + \varepsilon_p,
\end{equation*}
con $\varepsilon_p \in \mathbb{R} \sim E_p(\theta)$, de manera que $E_p$ es una variable aleatoria con vector de par\'ametros $\theta$, tal que $q_p(\varepsilon) = 0$. 

Es importante aclarar que $f_p(x) \in \mathbb{R}$ y $\theta$ son desconocidos. Asimismo, al igual que con la \textit{regresi\'on a la media}, se supone independencia entre los $\varepsilon$s, y, por lo tanto, hay independencia condicional entre las observaciones.

Otro aspecto importante a resaltar es que en este contexto la interpretaci\'on de $\varepsilon_p$ como el \textit{error aleatorio} ya no hace tanto sentido, y tendr\'ia que ser entendido m\'as como la dispersi\'on que siguen los datos alrededor de $f_p$.

\subsection{Modelo tradicional}

Cuando surgi\'o entre la comunidad estad\'istica el problema de \textit{regresi\'on sobre cuantiles}, inicialmente fue modelado bajo un enfoque no Bayesiano, como se describe en \cite{Yu_BayQuantReg}.  Posteriormente, Koenker \& Bassett (1978) retomaron esas ideas, y las aplicaron en el paradigma Bayesiano. 

Al igual que en la \textit{regresi\'on a la media}, el primer y m\'as popular modelo ha sido el lineal. Es decir, para $p \in (0,1)$ fijo arbitrario, se define
\begin{equation*}
    f_p(x) = x^T\beta_p, 
\end{equation*}
donde $\beta_p$ es el vector de coeficientes, dependiente de $p$.

\begin{defin}
    Se define a la funci\'on
    \begin{equation*}
        \rho_p(u) = u \times [pI_{(u>0)} - (1-p) I_{(u<0)})].
    \end{equation*}
    Se dice que una variable aleatoria u sigue una distribuci\'on asim\'etrica de Laplace ($u \sim AL_p(\sigma)$) si su funci\'on de densidad se escribe como
    \begin{equation*}
        w_p^{AL}(u|\sigma) = 
        \frac{p(1-p)}{\sigma}
        exp\left[
        -\rho_p
        \left(
        \frac{u}{\sigma}
        \right)
        \right],
    \end{equation*}
con $\sigma$ par\'ametro de escala.
\end{defin}

Data esta definici\'on, es posible darse cuenta que si $\varepsilon_p \sim AL_p(\sigma)$, entonces $q_p(\varepsilon_p) = 0$. Recordando que esta es la \'unica caracter\'istica necesaria para la variable de dispersi\'on, entonces es posible definir
\begin{equation*}
    \varepsilon_p \sim AL_p(\sigma).
\end{equation*}

El modelo, entonces, se puede reescribir como:
\begin{equation*}
    y | x, \beta_p, \sigma 
    \sim 
    AL_p(y - x^T\beta_p|\sigma).
\end{equation*}

Sea $\{(X,Y) | X \in \mathbb{R}^{m \times n}, Y \in \mathbb{R}^m\}$ el conjunto de datos observados. Por el Teorema de Bayes,
\begin{equation*}
\begin{aligned}
    \mathbb{P}(\beta_p,\sigma | Y, X) 
    &\propto \mathbb{P}(Y| X, \beta_p, \sigma) \times \mathbb{P}(\beta_p, \sigma), \\
\end{aligned}
\end{equation*}
donde $\mathbb{P}(Y| X, \beta_p, \sigma)$ es la verosimilitud de los datos observados y, debido a la independencia condicional, se puede calcular como 
\begin{equation*}
    \mathbb{P}(Y| X, \beta_p, \sigma)
    =
    \prod_{i=1}^m AL_p(y_i - x_i^T\beta_p|\sigma).
\end{equation*}

Por otro lado, $\mathbb{P}(\beta_p,\sigma^2)$ es la distribuci\'on inicial de los par\'ametros, para los que normalmente se usa
\begin{equation*}
    \beta_p,\sigma \sim \mathcal{NGI}(M,V,a,b). 
\end{equation*}

A diferencia del modelo tradicional de la \textit{regresi\'on a la media}, este modelo no es conjugado. Por lo tanto se requieren m\'etodos computacionales (como los que ser\'an descritos en el cap\'itulo 5) para aproximar la distribuci\'on posterior.

En el caso de la predicci\'on, si se tiene una nueva matriz de covariables $X_* \in \mathbb{R}^{r \times n}$, la inferencia con los datos observados se realiza de la siguiente manera:
\begin{equation*}
\begin{aligned}
    \mathbb{P}(Y_*|X_*,Y,X)
    &= \int \int \mathbb{P}(Y_*|X_*,\beta_p,\sigma) \times \mathbb{P}(\beta_p,\sigma|Y,X) d\sigma d\beta_p \\
    &= \int \int \prod_{i=1}^r AL_p(y_{i*} - x_{i*}^T\beta_p|\sigma) \times \mathbb{P}(\beta_p,\sigma|Y,X) d\sigma d\beta_p,
\end{aligned}
\end{equation*}
que tampoco tiene soluci\'on anal\'itica.

Si bien este modelo representa un gran avance, a\'un queda la posibilidad de retomar estas ideas y crear modelos m\'as flexibles, que capturen con mayor precisi\'on las particularidades de cada fen\'omeno y la interacci\'on entre las variables de salida y las covariables. En el siguiente cap\'itulo se discutir\'a la importancia de capturar mayor complejidad en la distribuciones, tanto de $f_p$, como de $\varepsilon_p$, mediante el de uso de m\'etodos no param\'etricos.
\input{Chapters/IV-NonParametric}
\chapter[Modelo GPDP para regresi\'on sobre cuantiles]{Modelo GPDP para regresi\'on sobre cuantiles}

\section{Definici\'on}

Despu\'es de analizar la introducci\'on de componentes no param\'etricos en las distribuciones, tanto de $f_p$, como de $\varepsilon_p$, a continuaci\'on se enunciar\'a el modelo central de esta tesis, con sus especificaciones correspondientes.

A partir de este punto, a dicho modelo se le denominar\'a \textbf{Modelo GPDP} (por las siglas en ingl\'es de procesos Gaussianos y procesos de Dirichlet).

Sea $\{(y_i,x_i)|i=1,...,m\}$ el conjunto de observaciones de la variable de respuesta y sus respectivas covariables, cuya relaci\'on se supone como
\begin{equation*}
    y = f_p(x) + {\varepsilon_p},
\end{equation*}
donde $f_p: \mathbb{R}^n \times \mathbb{R}$ es la funci\'on cuantil y ${\varepsilon_p} \in \mathbb{R}$ es el error aleatorio, ambos desconocidos.

Para reflejar la incertidumbre y el conocimiento previo del modelador, se supone a $f_p \sim \mathcal{GP}(m,k)$, con funci\'on de medias $m$ dada por el modelador y funci\'on de covarianza $k$ del tipo 2-\textit{exponencial}, con par\'ametro de rango fijo $\tau = 1$. Es decir,
\begin{equation*}
    k(x_i, x_j|\lambda) = \lambda \text{ } exp\{-\norm{x_i - x_j}_2\},
\end{equation*}
con $\lambda \sim GI(c_\lambda,d_\lambda)$, siendo $c_\lambda$ y $d_\lambda$ los par\'ametros de forma y escala, respectivamente, de una \textit{Gamma-Inversa}, mismos que deber\'an ser elegidos por el modelador. 

La raz\'on para fijar $\tau = 1$ es para simplificar el proceso computacional de inferencia que se ver\'a en la siguiente secci\'on, pero bien podr\'ia tambi\'en tener una distribuci\'on inicial que refleje la incertidumbre acerca de su valor.

En cuanto a la distribuci\'on inicial de $\varepsilon_p$, se supondr\'a un modelo de mezclas infinitas de Dirichlet, cuya distribuci\'on media $H$ del proceso de Dirichlet ser\'a una \textit{Gamma-Inversa}, con par\'ametros de forma $c_{DP}$ y escala $d_{DP}$, elegidos por el investigador.

En resumen, el Modelo GPDP queda descrito de la siguiente forma:
\begin{equation*}
\begin{aligned}
    y_i| f_p(x_i), z_i, \sigma_k^* &\sim AL_p({\varepsilon_p}_i = y_i - f_p(x_i) | \sigma_{z_i}), \\
    f_p|m, k, \lambda &\sim \mathcal{GP}(m,k(\lambda)|\lambda), \\
    \lambda &\sim GI(c_\lambda,d_\lambda), \\
    z_i | \pi &\sim Mult_\infty(\pi), \\
    \pi | \alpha &\sim GEM(\alpha), \\
    \sigma_k^* | c_{DP}, d_{DP} &\sim GI(\sigma_k|c_{DP}, d_{DP}),\\
    k(x_i, x_j | \lambda) &= \lambda \text{ } exp\{-\norm{x_i - x_j}_2\}.
\end{aligned}
\end{equation*}

\section{Inferencia con el simulador de Gibbs}

Dado que el modelo descrito no es conjugado, las distribuciones posteriores tienen que ser aproximadas mediante m\'etodos computacionales. Para hacer esto, se puede hacer uso de algoritmos MCMC (Markov chain Monte Carlo), y particularmente del simulador de Gibbs. \footnote{En caso de que el lector no est\'e familiarizado con este tipo de algoritmos, puede consultar una breve descripci\'on de ellos en el \autoref{chap:MCMC}.}

En este orden de ideas, a continuaci\'on se detallan las distribuciones condicionales posteriores de los par\'ametros del modelo, as\'i como la inclusi\'on de algunas variables latentes para permitir el funcionamiento del algoritmo. Es oportuno recordar que dichas distribuciones posteriores resultan de multiplicar la verosimilitud por la probabilidad inicial, como se revis\'o en el cap\'itulo 2 de este trabajo.

Antes de correr los algoritmos, usualmente resulta conveniente estandarizar los datos. En primer lugar, para que la estructura de covarianza tenga más sentido, ya que la escala de las covariables afectaría la correlación que existe entre los datos, al depender esta de la distancia entre ellas. Además, estandarizar los datos suele mejorar el rendimiento computacional de este tipo de algoritmos. Asimismo, vuelve m\'as sencillo definir el valor inicial de los hiper-par\'ametros, como se detallar\'a m\'as adelante.

\subsection{Actualizaci\'on del error}

Recordando que los centros de masa y los pesos de un proceso de Dirichlet son independientes, pueden ser actualizados por separado, con el inconveniente de que hay un n\'umero infinito de par\'ametros que actualizar. Para resolverlo, se utilizará el algoritmo de truncamiento del \textit{slice sampling}, propuesto por \cite{Kalli_Slice}, y adaptado para el modelo propuesto en esta tesis. A grandes rasgos consiste en truncar las posibles subpoblaciones a un número finito, el cual se actualizar\'a de forma din\'amica, de acuerdo a lo que vaya aprendiendo de los datos. 

Sea $\xi_1,\xi_2,\xi_3,...$ una secuencia positiva, generalmente elegida de forma determinista y decreciente. Sea $N$ una variable aleatoria con soporte en los n\'umeros naturales, una variable auxiliar incorporada al modelo.

\subsubsection{Actualizaci\'on de los centros de masa}

Para cada $k \in \{1,2,...,N\}$, se obtiene que 
\begin{equation*}
\begin{gathered}
    \sigma_k | \{{\varepsilon_p}_i, z_i | z_i = k\}, c, d \sim GI(\bar{c}_{DP}, \bar{d}_{DP}),\\
    \bar{c}_{DP} = c_{DP} + |\{i| z_i = k\}|, \\
    \bar{d}_{DP} = d_{DP} 
    + p \left[\sum_{\{i| z_i = k,\text{ }{\varepsilon_p}_i \geq 0\}} {\varepsilon_p}_i\right]
    + (1-p) \left[\sum_{\{i| \text{ } z_i = k,\text{ }{\varepsilon_p}_i < 0\}}  -{\varepsilon_p}_i\right].
\end{gathered}
\end{equation*}

\subsubsection{Actualizaci\'on de los pesos}

Sea $\bar{\pi}_k = \beta_k \prod_{j=1}^{k-1}(1 - \beta_j)$, de modo que para cada $k \in \{1,2,...,N\}$, la distribuci\'on condicional posterior de $\beta_k$ es
\begin{equation*}
\begin{aligned}
    \beta_k|\{z_i\}, a,b &\sim Beta(\bar{a}, \bar{b}), \\
    \bar{a} &= 1 + |\{i|z_i = k\}|, \\
    \bar{b} &= \alpha + |\{i|z_i > k\}|.
\end{aligned}
\end{equation*}

Dado que existe un n\'umero finito de posibles subpoblaciones, ya no se sigue propiamente la distribuci\'on $GEM$, sino un truncamiento de ella hasta la $N$-\'esima subpoblaci\'on. Posteriormente se realiza un reescalamiento de las probabilidades para que sumen 1. Es decir, se calcula
\begin{equation*}
\begin{aligned}
    \pi_k = \frac{\bar{\pi_k}}{\sum_{j=1}^N \bar{\pi_j}}
\end{aligned}
\end{equation*}

\subsubsection{Actualizaci\'on de las clases y variables de truncamiento}

Siguiendo el algoritmo de \cite{Kalli_Slice}, para cada observaci\'on $i \in \{1,...,m\}$, se obtiene
\begin{equation*}
\begin{aligned}
   u_i \sim U(0, \xi_{z_i}),
\end{aligned}
\end{equation*}
valor que se utiliza para actualizar la probabilidad de pertenencia a cada clase de la siguiente forma. Para cada $k \in \{1,2,...,N\}$,
\begin{equation*}
\begin{aligned}
   P(z_i = k| {\varepsilon_p}_i, \pi_k, \sigma_k)
   \propto
   \mathds{1}(u_i < \xi_k)
   \cdot
   \frac{\pi_k}{\xi_k}
   \cdot
   AL_p({\varepsilon_p}_i | \sigma_k).
\end{aligned}
\end{equation*}

Posteriormente se actualiza
\begin{equation*}
\begin{aligned}
   N = \max\{
    N_i|N_i=\max\{j|\xi_j > u_i\}, 
    i \in \{1,...,m\}
   \}.
\end{aligned}
\end{equation*}

\subsection{Actualizaci\'on de la tendencia}

Se define la siguiente variable aleatoria auxiliar, con la finalidad de anticipar si $\varepsilon_p = y - f_p(x)$ ser\'a positiva o negativa, y as\'i simplificar el c\'alculo de la actualizaci\'on de $f_p$.
\begin{equation*}
\begin{aligned}
    b_i | p, \sigma_i &\sim 
    \begin{cases}
        \frac{p}{\sigma_i} &prob = P({\varepsilon_p}_i \geq 0) = 1-p\\
        -\frac{1-p}{\sigma_i} &prob = P({\varepsilon_p}_i < 0) = p
    \end{cases},\\
\end{aligned}
\end{equation*}
de forma que $b = [b_1,...,b_m]^T$. 

\subsubsection{Actualizaci\'on de $\bm{f_p(X)}$}

Es pertinente recordar que la funci\'on de densidad de una observaci\'on $y_i$, debido a que sigue la distribuci\'on asim\'etrica de Laplace, se escribe
\begin{equation*}
\begin{aligned}
    P(y_i | f_p(x_i),\sigma_i) = 
    \frac{p(1-p)}{\sigma_i}
    exp \left\{-\rho_p
        \left(
            \frac{y_i-f_p(x_i)}{\sigma_i}
        \right)
    \right\}
    \mathds{1}_{(-\infty,\infty)}.
\end{aligned}
\end{equation*}

Una vez calculada la variable auxiliar $b$ que reci\'en se acaba de definir, dicha densidad se puede expresar de forma condicional como
\begin{gather*}
    P(y_i | f_p(x_i),\sigma_i,b_i) \propto
    \begin{cases}
        exp \left\{
            \frac{-p(y_i-f_p(x_i))}{\sigma_i}
        \right\} \mathds{1}_{\{y_i - f(x_i) \geq 0\}} 
        &\text{si } b_i > 0\\
        exp \left\{
            \frac{(1-p)(y_i-f_p(x_i))}{\sigma_i}
        \right\} \mathds{1}_{\{y_i - f(x_i) < 0\}} 
        &\text{si } b_i < 0\\
    \end{cases}.
\end{gather*}
Por lo tanto, la verosimilitud de las observaciones se puede calcular como
\begin{gather*}
    P(y | f_p(X),\sigma,b) \\
    \propto exp \left\{
        -\sum_{\{i|b_i > 0\}} \frac{p}{\sigma_i}(y_i-f_p(x_i)) -
        \sum_{\{i|b_i < 0\}} -\frac{1-p}{\sigma_i}(y_i-f_p(x_i))
    \right\} \\
    \prod_{\{i|b_i > 0\}}\mathds{1}_{\{y_i \geq f(x_i)\}}
    \prod_{\{i|b_i < 0\}}\mathds{1}_{\{y_i < f(x_i)\}}\\
    = exp \left\{
        -\sum_{\{i|b_i > 0\}} b_i(y_i-f_p(x_i)) -
        \sum_{\{i|b_i < 0\}} b_i(y_i-f_p(x_i))
    \right\} \\
    \prod_{\{i|b_i > 0\}}\mathds{1}_{\{y_i \geq f(x_i)\}}
    \prod_{\{i|b_i < 0\}}\mathds{1}_{\{y_i < f(x_i)\}}\\
    = exp \left\{-b^T(y-f_p(X))
    \right\}
    \prod_{\{i|b_i > 0\}}\mathds{1}_{\{y_i \geq f(x_i)\}}
    \prod_{\{i|b_i < 0\}}\mathds{1}_{\{y_i < f(x_i)\}}.\\
\end{gather*}
En esta peculiar verosimilitud, cada $f_p(x_i)$ estar\'a condicionada de manera excluyente a estar por arriba o por abajo de $y_i$. Por ese motivo, al multiplicar la verosimilitud por la distribuci\'on inicial Gaussiana de $f_p(X)$, se obtendr\'a como distribuci\'on posterior una Normal Truncada, misma que se detalla a continuaci\'on.
\begin{equation*}
\begin{aligned}
   f_p(X)|Y,X,M,b,\lambda &\sim TruncNormal(\bar{M}(X,b), K(X,X|\lambda), \gamma, \eta), \\
   \bar{M}(X,b) &= M(X) + K(X,X|\lambda)b, \\
   \gamma_i &= 
   \begin{cases}
    -\infty & \text{si }b_i > 0 \\
    y_i & \text{si }b_i < 0
   \end{cases},\\
   \eta_i &= 
   \begin{cases}
    y_i & \text{si }b_i > 0 \\
    \infty & \text{si }b_i < 0
   \end{cases},
\end{aligned}
\end{equation*}
donde $\gamma$ es el vector de l\'imites inferiores y $\eta$ es el vector de l\'imites superiores de la distribuci\'on Normal truncada.

Debido a que $f_p(X)$ se encuentra en la verosimilitud  \'unicamente como un elemento de primer orden, al hacer la multiplicaci\'on con la distribuci\'on inicial Normal, la varianza queda exactamente igual. La actualizaci\'on se da \'unicamente en la media, como una perturbaci\'on de la media inicial, dada por las covarianzas que tiene cada observaci\'on respecto a las dem\'as, as\'i como el signo de $b_i$ para cada una.

\subsubsection{Actualizaci\'on del par\'ametro de escala}

Condicional a los dem\'as valores obtenidos, se obtiene la distribuci\'on posterior  de $\lambda$ como
\begin{equation*}
\begin{gathered}
   P(\lambda|X,M(X),f_p(X),b,c_\lambda,d_\lambda) 
   \propto
   \lambda^{-\bar{c}_\lambda-1}
   \cdot
   exp\left\{- \frac{\bar{d}_\lambda}{\lambda}\right\}
   \cdot
   exp\left\{-\bar{B} \lambda\right\}, \\
   \bar{c}_\lambda = c_\lambda + \frac{p}{2}, \\
   \bar{d}_\lambda = d_\lambda + \bar{F}, \\
   \bar{F} = \frac{1}{2}(f_p(X)-M(X))^T [K(X,X|\lambda=1)^{-1}] (f_p(X)-M(X)), \\
   \bar{B} = \frac{1}{2}b^T [K(X,X|\lambda=1)] b.
\end{gathered}
\end{equation*}

\section{Predicci\'on}

Una de las desventajas de los modelos no param\'etricos es que, a diferencia de los modelos param\'etricos, es complicado interpretar los resultados del ajuste del modelo. Por ello, resulta particularmente importante la faceta de la predicci\'on, que es donde m\'as se puede explotar la flexibilidad de los modelos no param\'etricos, debido a que el objetivo es tener precisi\'on en la estimaci\'on. Espec\'ificamente esta secci\'on se enfocar\'a en la predicci\'on de $f_p$, que es el par\'ametro de mayor inter\'es del modelo.

Debido al uso del simulador de Gibbs, despu\'es de realizar el ajuste se cuenta con un conjunto grande de realizaciones aproximadas de $f_p(X)$, provenientes de las cadenas de Markov.

Recordando lo visto en la secci\'on 4.2.4, cuando se tienen valores de $f_p(X)$, es posible usar la propiedad de la \textit{Normal condicional} para realizar predicci\'on. Sea $X \in \mathbb{R}^m \times \mathbb{R}^n$ la matriz de datos originales, $X_* \in \mathbb{R}^r \times \mathbb{R}^n$ la matriz de datos a predecir, $f_p(X)$ una realizaci\'on de la distribuci\'on posterior correspondiente a X, y $f_p(X_*)$ el vector aleatorio de los datos a predecir. Se tiene entonces que 
\begin{equation*}
    f_p(X_*)|f_p(X) 
    \sim \mathcal{N}
    (\bar{M}(X,X_*),\bar{K}(X,X_*|\lambda)),
\end{equation*}
con
\begin{equation*}
\begin{aligned}
    \bar{M}(X,X_*) &= M(X_*) + K(X_*,X)K(X,X)^{-1}(f_p(X) - M(X)), \\
    \bar{K}(X,X_*|\lambda) &= 
    \lambda
    \times
    \left[
    K(X_*,X_*) -
    K(X_*,X)K(X,X)^{-1}K(X,X_*)
    \right]
    .
\end{aligned}
\end{equation*}
donde $K(X_1,X_2) = K(X_1,X_2|\lambda=1)$, y $X_1$ y $X_2$ pueden ser $X$ o $X_*$.

Por lo antes descrito, es posible obtener una realizaci\'on de $f_p(X_*)$ simulando de dicha distribuci\'on Normal. De esta manera, por cada valor de $f_p(X)$ y $\lambda$ en la cadena de Markov, se simula una realizaci\'on de $f_p(X_*)$, y entonces es posible aproximar la distribuci\'on posterior de $q_p(y|x)$, para los datos $X_*$.

\section{Hiper-par\'ametros iniciales del modelo}

Una complicaci\'on que puede tener un modelo jer\'arquico, como el GPDP, es que es complicado darle una interpretaci\'on intuitiva a los hiper-par\'ametros de las diversas capas, de forma que el conocimiento previo del modelador pueda reflejarse en valores asignados a ellos.

Para mitigar este problema, a continuaci\'on se proponen una serie de heur\'isticas para definirlos, mismas que se derivan de algunas ideas que me parecen sensatas, pero no se originan de ning\'un cuerpo axiom\'atico y bien podr\'ian ser mejoradas. Tambi\'en es importante aclarar que por lo comentado al inicio de la secci\'on 5.2, para todas ellas se pensar\'a que los datos est\'an estandarizados.

\subsection{Funci\'on de medias $m$}

Este es el hiper-par\'ametro al que se le puede dar una mayor interpretaci\'on, debido a que representa el nivel donde el modelador estima que se encontrar\'a el cuantil p-\'esimo de $y$, para cada valor de las covariables $x$.

Para simplificar el proceso de definici\'on de la funci\'on de medias $m$, se puede partir de la hip\'otesis que es constante, y, por lo tanto, las variaciones son \'unicamente producto de la varianza de $\varepsilon_p$. Dada la estructura de probabilidad posterior, la media de $f_p$ podr\'a actualizarse si los datos brindan informaci\'on suficiente para suponer lo contrario. 

Una vez aceptada esta estructura, resta asignar el valor constante que tomar\'a. Si el modelador tiene una idea del nivel donde espera los datos, puede asignar la constante $c$. En caso de no tenerla, un valor que se suele usar en diversos contextos es el de $c=0$, de forma que
\begin{equation*}
\begin{aligned}
    m:\mathbb{R}^n &\rightarrow \mathbb{R}, \text{tal que }\\
    m&(x) = c.
\end{aligned}
\end{equation*}

\subsection{\textit{Gamma-Inversa}s de $\lambda$ y el Proceso de Dirichlet}

Tanto $c_\lambda$ y $d_\lambda$, como $c_{DP}$ y $d_{DP}$ son par\'ametros de distribuciones \textit{Gamma-Inversa}. Es oportuno recordar que si $U \sim \mathcal{GI}(c,d)$, entonces
\begin{equation*}
\begin{aligned}
    \mathbb{E}[U] &= \frac{d}{c-1}, \text{ } c>1,\\
    Var(U) &= \frac{d^2}{(c-1)^2(c-2)}, \text{ } c>2.
\end{aligned}
\end{equation*}

En este orden de ideas, si se elige $c = 2$, $Var(U)$ ser\'a infinita y $\mathbb{E}[U] = d$. Asignar a $c_\lambda$ y $c_{DP}$ de esta manera permitir\'a darle a $d_\lambda$ y $d_{DP}$ el valor que se piense como el mejor estimador puntual \textit{a priori} de $\lambda$ y $\sigma$, pero con una varianza tal que permitir\'a a los datos tener el peso principal en la actualizaci\'on del modelo. 

Debido a la estandarizaci\'on de los datos, la varianza muestral de $y$ es igual a 1. Es posible pensarla como el resultado de sumar la varianza de $f_p(x)$ y la de $\varepsilon_p$, que adem\'as se suponen independientes. Entonces, se puede definir una heur\'istica tal que $Var(f_p(x)) = \frac{1}{2}$ y $Var(\varepsilon_p) = \frac{1}{2}$, a falta de mayor informaci\'on.

La varianza de $f_p(x)$ es igual a $\lambda$, por lo que lo coherente con lo dicho en los p\'arrafos anteriores ser\'a asignar $d_\lambda = \frac{1}{2}$. 

Por el otro lado, si \'unicamente para este ejercicio, y con el af\'an de volver an\'alitico el c\'alculo, se piensa a $\varepsilon_p \sim AL_p(\sigma = d_{DP})$. Entonces, su varianza estar\'ia dada por
\begin{equation*}
    Var(\varepsilon_p) = 
    \left[\frac{d_{DP}}{p(1-p)}\right]^2
    (1-2p(1-p)).
\end{equation*}
Dado que se fijar\'a $Var(\varepsilon_p) = \frac{1}{2}$, por la heur\'istica antes mencionada, despejando es posible obtener que
\begin{equation*}
    d_{DP} = \frac{p(1-p)}{\sqrt{2(1-2p(1-p))}}.
\end{equation*}

\subsection{Par\'ametro de concentraci\'on $\alpha$}

Este es el par\'ametro m\'as dif\'icil de definir, por su complejidad de interpretaci\'on. Pero cabe recordar que el valor de $\alpha$ tiene una relaci\'on positiva con el n\'umero de subpoblaciones. 

De hecho, sea $\bar{m}$ el n\'umero de subpoblaciones y $m$ el n\'umero de datos de entrenamiento, \cite{Yee_DirProc} expone que
\begin{equation*}
    \mathbb{E}[\bar{m}|\alpha, m] 
    \simeq 
    \alpha
    \log 
    \left(
        1 + \frac{m}{\alpha}
    \right)
    \text{, para } m, \alpha \gg 0.
\end{equation*}

Si se define $\alpha = \frac{\sqrt{m}}{2}$, se tiene que
\begin{equation*}
\begin{aligned}
    \mathbb{E}[\bar{m}|m] 
    &\simeq 
    \frac{\sqrt{m}}{2}
    \times
    \log 
    \left(
        1 + 2\sqrt{m}
    \right)\\
    &\simeq
    \frac{m}{7} 
    \text{, para } m \approx 100.
\end{aligned}
\end{equation*}

Es decir, si se tienen alrededor de 100 observaciones, el n\'umero esperado de subpoblaciones ser\'a alrededor de la s\'eptima parte de las observaciones. Valor que a falta de mayor exploraci\'on en este tema, no suena descabellado. De nuevo, a mayor cantidad de datos, se tendr\'a una menor dependencia de esta arbitraria decisi\'on inicial.

\section{Consideraciones sobre la bondad de ajuste}

El modelo GPDP descrito en este trabajo pretende ser una opci\'on m\'as de modelo de regresi\'on, particularmente \'util cuando es inter\'es del modelador aproximar alg\'un cuantil en espec\'ifico de la variable de respuesta, dados los valores de las variables explicativas.

Sin embargo, no est\'a de m\'as recordar que existen muchos otros modelos de regresi\'on, a la media y sobre cuantiles, lineales y no lineales, param\'etricos y no param\'etricos. Todos ellos cumplen el mismo cometido de estimar la distribuci\'on condicional de $y|x$.

En la siguiente secci\'on se har\'a un an\'alisis de los resultados obtenidos por el modelo GPDP, que har\'a uso de herramientas b\'asicas, como la exploraci\'on visual o la correlaci\'on. En este sentido, hay una evidente \'area de mejora, ya que se podr\'ia encontrar alguna medida robusta, estad\'isticamente hablando, de bondad de ajuste. \'Esta permitir\'ia saber qu\'e tan bueno resulta el modelo para describir un conjunto dado de datos, y permitir\'ia compararlo con otros, para seleccionar cu\'al es el que presenta mejores resultados. 

Desafortunadamente, realizar esto con un modelo con la estructura del GPDP resulta complejo y la literatura para medir la bondad de ajuste de modelos de regresi\'on bayesianos, sobre cuantiles y no parame\'etricos, est\'a apenas en sus pininos. Por lo tanto, una estimaci\'on robusta de la bondad de ajuste del modelo GPDP y la selecci\'on del mejor modelo quedan fuera del alcance de este trabajo.


\section{Paquete \textit{GPDPQuantReg} en R}

Todas las ideas expuestas en este cap\'itulo han sido implementadas en el paquete \textit{GPDPQuantReg} del lenguaje de programaci\'on R, mismo que puede ser encontrado en el repositorio de Github\faGithub \space titulado: \textbf{opardo/GPDPQuantReg}.

Al momento de escribir este trabajo, cuenta con tres funciones p\'ublicas: \textit{GPDPQuantReg}, para ajustar el modelo con el simulador de Gibbs; \textit{predict}, para realizar predicci\'on en un nuevo conjunto de datos del modelo ajustado; y \textit{diagnose}, para realizar el diagn\'ostico de la ergodicidad, la autocorrelaci\'on, la correlaci\'on cruzada y la traza de las cadenas de Markov, para los distintos par\'ametros.

A continuaci\'on se expone un ejemplo de uso, el cual es similar a lo que se realiz\'o para obtener los resultados del cap\'itulo siguiente.

\lstinputlisting{R/package_example.R}

\newpage
\chapter[Aplicaciones]{Aplicaciones}

\section{Metodolog\'ia de simulaci\'on de datos y ajuste de los modelos}

Para la elaboraci\'on de este cap\'itulo se ajustaron modelos de diversos cuantiles, a distintos conjuntos de datos. Los datos se obtuvieron de la siguiente manera. Sea $y \in \mathbb{R}$ el valor de la variable de respuesta, $x \in \mathbb{R}$ su respectiva covariable, $g: \mathbb{R} \rightarrow \mathbb{R}$ una funci\'on denominada \textit{tendencia} y $\omega(x) \in \mathbb{R}$ un error aleatorio, se simul\'o
\begin{equation*}
    y = g(x) + \omega(x).
\end{equation*}
Cabe resaltar que se utiliza la notaci\'on $g$, y $\omega$, para diferenciarlas de $f_p$ y $\varepsilon_p$, que son las usadas en el planteamiento de los modelos te\'oricos, vistos en los cap\'itulos anteriores. La raz\'on es que no necesariamente se cumple que $q_p(\omega|x) = 0$. De hecho, como se ver\'a en las siguientes l\'ineas, la funci\'on que se querr\'a estimar no ser\'a $g$, sino las correspondientes $f_p$ para diferentes cuantiles de inter\'es.

La funci\'on real del cuantil $p$\textit{-\'esimo} de $y|x$ se puede obtener como
\begin{equation*}
    q_p(y|x) = g(x) + q_p(\omega|x),
\end{equation*}
funci\'on que en la construcci\'on te\'orica se denomin\'o como $f_p$, y la cual se desea estimar para diversos valores de $p$.

En todos los casos presentados a continuaci\'on se ajust\'o el modelo para el cuantil $0.5$\textit{-\'esimo}, por ser la mediana y una medida de tendencia central alternativa a la media. Tambi\'en se model\'o el cuantil $0.95$\textit{-\'esimo}, ya que es un valor extremo. Finalmente, se busco describir al cuantil $0.25$\textit{-\'esimo}, debido a que el primer cuartil es de inter\'es dentro de algunos contextos, y no es medida de tendencia central, pero tampoco un valor extremo.

Para todos los casos se simularon 60 datos sin reemplazo dentro del intervalo $(-15,15)$, con un refinamiento de 2 decimales. Las razones de elegir esa cantidad de datos, que podr\'ia parecer reducida, fueron varias. Primero, para validar el poder predictivo del algoritmo y su capacidad de ajuste a los valores originales, aún cuando se tienen pocos datos. Segundo, para obtener viabilidad en tiempos de ejecuci\'on, ya que el algoritmo simula en cada iteraci\'on, entre otras variables aleatorias, normales multivariadas y truncadas, con dimensi\'on igual al n\'umero de datos. Dicha simulaci\'on utiliza el m\'etodo de aceptaci\'on y rechazo, por lo que entre mayor sea la dimensi\'on, usualmente requiere de mayor esfuerzo computacional, y, por ende, el tiempo total del ajuste del modelo aumenta de manera significativa. Finalmente, la tercer raz\'on es que esa cantidad de datos permitir\'a observar qu\'e influencia puede tener un dato particular en el ajuste general del modelo.

Una vez simulados los datos, se ajustaron el modelo tradicional de regresi\'on a la media (descrito en la secci\'on \ref{trad_mean_reg} de este trabajo) y el GPDP. Es importante recordar que el primero, adem\'as de ser sobre la media, es lineal en los  par\'ametros de las covariables y con error normal. Por otro lado, el modelo GPDP propone ideas alternativas, como centrar su estimaci\'on en el cuantil de inter\'es, y utilizar componentes no param\'etricos, tanto para modelar la tendencia, como el error. En medio queda una amplia gama de modelos, como el modelo tradicional de regresi\'on sobre cuantiles, o alg\'un modelo de regresi\'on a la media que contemple errores no param\'etricos, por ejemplo. 

Sin embargo, decidi\'i realizar esa particular comparaci\'on para descubrir qu\'e tan bueno podr\'ia ser un modelo que contempla ideas alternativas en todas las facetas de la regresi\'on (estad\'istico a modelar, tendencia y error aleatorio), cuando se le comparara con el que utiliza \'unicamente supuestos tradicionales. Eso no significa que si se mezclan algunas ideas alternativas, con algunas tradicionales, no se pueda obtener un tercer modelo, que resulte a\'un mejor. Pero el alcance de esta tesis se limita a comparar \'unicamente esos dos modelos.

Para el ajuste del modelo tradicional se utiliz\'o un poliniomio de orden 5. Es decir, se utilizaron como covariables los valores de $x, x^2, \ldots, x^5$, adem\'as del t\'ermino independiente. A\'un as\'i, dicho ajuste fue simple, computacionalmente hablando, ya que requiri\'o \'unicamente de multiplicaciones matriciales para actualizar sus par\'ametros posteriores. Como par\'ametros iniciales de la distribuci\'on Normal de $\beta$, se utiliz\'o el vector de medias con valores iguales a 0, y como matriz de varianzas y covarianzas, la identidad. En el caso de la \textit{Gamma-Inversa} de $\sigma^2$, el valor de forma usado fue 2, y el de escala, 1.

Por otra parte, el ajuste del GPDP involucr\'o el c\'alculo de una cadena de Markov con ciertas particularidades. En todos los casos se desecharon las primeras $5,000$ simulaciones, para as\'i evitar dependencia de los valores iniciales. Posteriormente, se simularon $15,000$ valores para cada par\'ametro, pero de ellos \'unicamente se guard\'o uno de cada cinco, con la intenci\'on de buscar replicar el comportamiento independiente, que se supone entre los valores. Por lo tanto, se utilizaron $3,000$ valores simulados para aproximar la distribuci\'on posterior de cada par\'ametro. Dicho ajuste se realiz\'o utilizando el paquete \textit{GPDPQuantReg} en R, mismo que, como se detalló en el cap\'itulo anterior, implementa el modelo GPDP para la regresi\'on sobre cuantiles.

Posteriormente, se realiz\'o la predicci\'on para cada $x \in [-15,15]$, tal que $x$ tuviera valor cero para todos los decimales a partir de las cent\'esimas. Es decir, $-15.0,-14.9,\ldots,14.9,15.0$. Se utilizaron ambos modelos para simular $3,000$ predicciones de cada $f_p(x)$ \footnote{El algoritmo utilizado para simular cuantiles del modelo tradicional de regresi\'on a la media se detalla en el \autoref{trad_mean_alg}.}. Posteriormente, se calcularon los cuantiles  muestrales $0.025$ y $0.975$\textit{-\'esimos}, para crear el intervalo de probabilidad al $95\%$ de la estimaci\'on posterior de la predicci\'on. Adem\'as, se calcul\'o el cuantil $0.5$\textit{-\'esimo} (la mediana), para tener una medida de tendencia central de dicha estimaci\'on. 

Finalmente, se obtuvieron m\'etricas que permitieran comparar a grandes rasgos la calidad de estimaci\'on de ambos modelos. Se contrast\'o la estimaci\'on mediana de la predicci\'on con el valor real del cuantil, y se obtuvo el error cuadr\'atico medio, as\'i como la correlaci\'on al cuadrado. Adem\'as se calcul\'o qu\'e porcentaje de los valores reales cayeron dentro del intervalo de probabilidad al 95\%, estimado para la predicci\'on.

\section{Conjuntos de datos simulados}

A continuaci\'on se presentan los diversos conjuntos de datos simulados. Primero, se expone uno que cumple con los supuestos de la regresi\'on tradicional a la media. Posteriormente, se presentan otros que desaf\'ian alguno de sus supuestos, con el fin de comparar c\'omo lidian dicho modelo, y el modelo GPDP, con tal inconveniente. Finalmente, se comparan los resultados.

\subsection{Supuestos tradicionales de regresi\'on a la media}

El primer caso que se presenta es el de un conjunto de datos que se simul\'o con las condiciones com\'unmente supuestas dentro de los modelos de regresi\'on a la media: funci\'on de tendencia polinomial y error normal alrededor de ella, independiente del valor de $x$. Es decir, los errores presentaron homocedasticidad. Expresado en t\'erminos matem\'aticos,
\begin{equation*}
    g(x) = \frac{1}{1000}x^3 - \frac{1}{40}x^2 - \frac{1}{10}x + 6,
\end{equation*}
mientras el error se distribuy\'o $\omega(x) \sim \mathcal{N}(0,1)$. Los datos simulados se pueden observar en la figura \ref{sample_classic}.

\begin{figure}[H]
	\centering
	\caption{Datos simulados con los supuestos tradicionales de regresi\'on.}
	\includegraphics[width=0.7\textwidth]{Figures/Simulation/classic/sample.png}
	\label{sample_classic}
\end{figure}

Posteriormente se ajustaron los modelos y se realiz\'o predicci\'on de la forma en la que se explic\'o anteriormente, obteniendo los siguientes resultados.

\begin{table}[H]
\centering
\caption{Error cuadrático medio de datos que cumplen supuestos tradicionales de regresi\'on.} 
\begin{tabular}{ccc}
  \hline
Cuantil & Modelo Tradicional & Modelo GPDP \\ 
  \hline
0.95 & 0.84 & 0.83 \\ 
  0.50 & 0.02 & 0.19 \\ 
  0.25 & 0.23 & 0.16 \\ 
   \hline
\end{tabular}
\label{mse_classic}
\end{table}

\begin{table}[H]
\centering
\caption{Correlación al cuadrado de datos que cumplen supuestos tradicionales de regresi\'on.} 
\begin{tabular}{ccc}
  \hline
Cuantil & Modelo Tradicional & Modelo GPDP \\ 
  \hline
0.95 & 0.99 & 0.91 \\ 
  0.50 & 0.99 & 0.94 \\ 
  0.25 & 0.99 & 0.96 \\ 
   \hline
\end{tabular}
\label{corr_classic}
\end{table}

\begin{table}[H]
\centering
\caption{Porcentaje de valores reales dentro del intervalo de confianza de datos que cumplen supuestos tradicionales de regresi\'on.} 
\begin{tabular}{ccc}
  \hline
Cuantil & Modelo Tradicional & Modelo GPDP \\ 
  \hline
0.95 & 68\% & 96\% \\ 
  0.50 & 100\% & 100\% \\ 
  0.25 & 100\% & 100\% \\ 
   \hline
\end{tabular}
\label{within_classic}
\end{table}

\begin{figure}[htbp]
	\centering
	\caption{Ajuste de los modelos Tradicional y \textit{GPDP}, para un conjunto de datos que cumplen los supuestos tradicionales de regresi\'on.}
	\includegraphics[width=\textwidth]{Figures/Simulation/classic/predictions.png}
	\captionsetup{singlelinecheck=off,font=footnotesize}
    \caption*{Nota: La l\'inea roja representa el valor real de cada cuantil, la l\'inea azul representa la mediana de la distribuci\'on posterior predictiva y el \'area gris su intervalo de probabilidad al 95\%.}
	\label{models_classic}
\end{figure}

Como era de esperarse, el modelo tradicional replic\'o de forma muy atinada el comportamiento polinomial de los datos. La m\'etrica que mejor refleja esta situaci\'on es la correlaci\'on al cuadrado, en donde dicho modelo obtuvo pr\'acticamente un valor perfecto para los tres cuantiles ajustados. Sin embargo, si bien no alcanz\'o los mismos niveles, el modelo GPDP tambi\'en obtuvo valores altos.

De hecho, si observa el cuadro \ref{mse_classic}, el lector podr\'a notar que para el cuantil $0.95$\textit{-\'esimo}, el error cuadr\'atico medio es pr\'acticamente el mismo, y para el cuantil $0.25$\textit{-\'esimo} es menor. Es decir, mientras que el modelo tradicional comprendi\'o de mejor manera las subidas y bajadas puntuales de la funci\'on, el modelo GPDP fue m\'as certero para estimar el nivel, conforme los valores se alejaron de la mediana. 

Lo anterior es confirmado por la tabla \ref{within_classic}, ya que para el caso del ajuste tradicional del cuantil $0.95$\textit{-\'esimo} \'unicamente el $68\%$ de los valores reales del cuantil quedaron dentro del angosto intervalo de probabilidad. Mi hip\'otesis es que dicho error de estimaci\'on del nivel proviene de un error min\'usculo de ajuste de la distribuci\'on de $\sigma^2$, el cual queda totalmente desapercibido en la estimaci\'on casi perfecta de la mediana, pero en un valor extremo, como es el caso del cuantil $0.95$\textit{-\'esimo}, s\'i se ve reflejado.

Hay comentarios del ajuste de este conjunto de datos que coincidir\'an con otros, por lo que en la secci\'on \ref{models_comp} se har\'an apuntes generales acerca de ellos.

Por otro lado, como se detall\'o en el cap\'itulo anterior, con el uso del paquete \textit{GPDPQuantReg} tambi\'en es posible obtener diagn\'osticos de las cadenas de Markov del modelo GPDP, los cuales se explican de manera m\'as detallada en el \autoref{chap:MCMC}. Como ejemplo, se presentan los del ajuste del cu\'antil $0.5$-\textit{\'esimo} para los datos de esta subsecci\'on, en la figura \ref{diag_sgse}. Los par\'ametros que se analizan son la $\lambda$, par\'ametro de escala de la funci\'on covarianza del proceso Gaussiano; y, adem\'as, se eligi\'o una observaci\'on de forma aleatoria, para la que se dio seguimiento al valor de $\sigma$ y $f_p(x)$. 

\begin{figure}[H]
	\centering
	\caption{Diagn\'osticos de las cadenas de Markov del cu\'antil $0.5$-\textit{\'esimo}, de datos que cumplen los supuestos tradicionales de regresi\'on.}
	\subfloat[Ergodicidad]{
	    \includegraphics[width=0.45\textwidth]{Figures/Simulation/simple_g_simple_error/ergodicity.png}
	}
	\subfloat[Autocorrelaci\'on]{
	    \includegraphics[width=0.45\textwidth]{Figures/Simulation/simple_g_simple_error/autocorrelation.png}
	}\\
	\subfloat[Correlaci\'on cruzada]{
	    \includegraphics[width=0.45\textwidth]{Figures/Simulation/simple_g_simple_error/crosscorrelation.png}
	}
	\subfloat[Traza]{
	    \includegraphics[width=0.45\textwidth]{Figures/Simulation/simple_g_simple_error/trace.png}
	}
	\label{diag_sgse}
\end{figure}

La ergodicidad parece adecuada, debido a que después de las primeras iteraciones, el promedio del valor de los par\'ametros se estabiliza en un nivel relativamente fijo para todas las iteraciones siguientes. La autocorrelaci\'on es excelente, ya que para los tres par\'ametros validados se muestra el valor de 1 en la correlaci\'on consigo mismo (por construcci\'on, siempre es as\'i), pero a partir del siguiente valor en la cadena de Markov, la correlaci\'on es pr\'acticamente 0. Es decir, los valores simulados s\'i parecen haber sido simulados de la distribuci\'on posterior, de forma independiente entre s\'i.

La correlaci\'on cruzada entre par\'ametros tambi\'en es cercana a 0, por lo que la simulaci\'on de un par\'ametro de la cadena no muestra una gran dependencia del valor de otro par\'ametro, en la misma iteraci\'on. Por ello, hay motivos para pensar que la cadena de Markov no cay\'o en ciclos distintos a su distribuci\'on estacionaria. De hecho, esto se puede confirmar al observar la traza de $f$, ya que no muestra comportamientos c\'iclicos y aunque se centra en cierto intervalo de valores, tambi\'en hay simulaciones que salen de \'el de forma espor\'adica.

\subsection{Error de colas pesadas}

En el siguiente conjunto de datos se utiliz\'o la distribuci\'on Cauchy para simular el error. A pesar de ser sim\'etrica y centrada en 0, dicha distribuci\'on es peculiar debido a que su esperanza y varianza no est\'an definidas. El motivo de ello son sus colas pesadas, mismas que originan con cierta frecuencia valores muy alejados del centro. Adem\'as, se introdujo un componente sinoidal en la funci\'on tendencia, mismo que puede ser aproximado con un polinomio tanto como se quiera, utilizando t\'erminos de su serie de Taylor. Sin embargo, tambi\'en tiene un valor absoluto, que requiere un n\'umero grande de t\'erminos para poder ser aproximado certeramente por la familia de los polinomios.

La expresi\'on de la funci\'on tendencia utilizada es \begin{equation*}
    g(x) = \frac{1}{4}|x| + sen(x).
\end{equation*}
Por otro lado, el error se distribuy\'o $\omega(x) \sim Cauchy(0,0.1)$. Los datos obtenidos se pueden observar en la figura \ref{sample_heavy_tails}.

\begin{figure}[H]
	\centering
	\caption{Datos simulados con error de colas pesadas.}
	\includegraphics[width=0.7\textwidth]{Figures/Simulation/heavy_tails/sample.png}
	\label{sample_heavy_tails}
\end{figure}

Como se puede observar, la mayor\'ia de los datos est\'an concentrados alrededor de la funci\'on tendencia. Sin embargo, ocasionalmente los valores se disparan y se alejan ampliamente de dicha concentraci\'on. Claros ejemplos son el dato que est\'a muy por debajo de los dem\'as, cerca de $x = -10$, y uno que est\'a muy por arriba, alrededor de $x = -5$. Lo interesante de este ejemplo fue ver c\'omo se comportaban ambos modelos ante dichos datos at\'ipicos. Los resultados se muestran a continuaci\'on.

\begin{table}[H]
\centering
\caption{Error cuadrático medio de datos con error de colas pesadas.} 
\begin{tabular}{ccc}
  \hline
Cuantil & Modelo Tradicional & Modelo GPDP \\ 
  \hline
0.95 & 5.42 & 0.18 \\ 
  0.50 & 0.61 & 0.75 \\ 
  0.25 & 1.89 & 1.12 \\ 
   \hline
\end{tabular}
\label{mse_heavy_tails}
\end{table}

\begin{table}[H]
\centering
\caption{Correlación al cuadrado de datos con error de colas pesadas.} 
\begin{tabular}{ccc}
  \hline
Cuantil & Modelo Tradicional & Modelo GPDP \\ 
  \hline
0.95 & 0.64 & 0.91 \\ 
  0.50 & 0.64 & 0.64 \\ 
  0.25 & 0.64 & 0.57 \\ 
   \hline
\end{tabular}
\label{corr_heavy_tails}
\end{table}

\begin{table}[H]
\centering
\caption{Porcentaje de valores reales dentro del intervalo de confianza de datos con error de colas pesadas.} 
\begin{tabular}{ccc}
  \hline
Cuantil & Modelo Tradicional & Modelo GPDP \\ 
  \hline
0.95 & 9\% & 98\% \\ 
  0.50 & 57\% & 100\% \\ 
  0.25 & 41\% & 96\% \\ 
   \hline
\end{tabular}
\label{within_heavy_tails}
\end{table}

\begin{figure}[htbp]
	\centering
	\caption{Ajuste de los modelos Tradicional y \textit{GPDP}, para un conjunto de datos con error de colas pesadas.}
	\includegraphics[width=\textwidth]{Figures/Simulation/heavy_tails/predictions.png}
	\captionsetup{singlelinecheck=off,font=footnotesize}
    \caption*{Nota: La l\'inea roja representa el valor real de cada cuantil, la l\'inea azul representa la mediana de la distribuci\'on posterior predictiva y el \'area gris su intervalo de probabilidad al 95\%.}
	\label{models_heavy_tails}
\end{figure}

Una vez obtenidos los resultados de ambos modelos, resulta interesante observar varios aspectos de la figura \ref{models_heavy_tails}. Tocando el tema de la influencia de valores at\'ipicos, parece claro que el modelo tradicional los manej\'o estimando valores muy altos para $\sigma^2$, del error que supone normal. Esto ocasion\'o que para los cuantiles distintos a la mediana, estimara que estaban m\'as lejos del centro, de lo que realmente estaban. Dicha situaci\'on se puede confirmar utilizando el cuadro \ref{within_heavy_tails}, ya que s\'olo el $41\%$ de los datos reales cay\'o dentro del intervalo de probabilidad estimado para el cuantil $0.25$\textit{-\'esimo}, y peor a\'un para el $0.95$\textit{-\'esimo}, donde s\'olo el $9\%$ lo hizo.

En cuanto a la estimaci\'on de la tendencia, el modelo tradicional interpret\'o que los datos ven\'ian de lo que parece ser una par\'abola, con m\'inimo y centro en 0, acerc\'andose relativamente al comportamiento del valor absoluto. Desafortunadamente no logr\'o atrapar el comportamiento sinoidal y atribuy\'o las fluctuaciones al error aleatorio. Eso se puede ver claramente en su estimaci\'on del cuantil $0.5$\textit{-\'esimo}, donde captur\'o el nivel de manera adecuada, pero como muestra la figura \ref{models_heavy_tails} y el cuadro \ref{corr_heavy_tails}, muchos valores salieron tanto por arriba, como por abajo, del intervalo.

Por otro lado, el modelo GPDP hizo una estimaci\'on muy certera en la mayor\'ia de los puntos, como se refleja en los cuadros \ref{mse_heavy_tails} y \ref{within_heavy_tails}. La vecindad alrededor de $x = 10$, que coincide con el valor at\'ipico m\'as dram\'atico, es donde el ajuste pudo haber sido mucho mejor. Sin embargo, creo que es importante recalcar que el problema se limit\'o a un error de estimaci\'on local, que no afect\'o a los puntos m\'as alla de cierta pequeña vecindad. De hecho, el siguiente valor at\'ipico en magnitud, alrededor de $x = 5$, fue manejado aceptablemente por el modelo para los tres cuantiles.

\subsection{Heterocedasticidad}

Otro supuesto com\'un en muchos modelos de regresi\'on es la homocedasticidad, es decir, que los errores no dependen de $x$. Para este ejemplo particular decid\'i observar c\'omo afectar\'ia a la estimaci\'on el hecho de que la distribuci\'on de los errores aleatorios dependiera del valor de la variable explicativa, es decir, datos que presentaran heterocedasticidad.

Para ello se us\'o como funci\'on de tendencia al polinomio de primer orden 
\begin{equation*}
    g(x) = \frac{1}{10}x + 2,
\end{equation*}
y la complejidad de estimaci\'on recay\'o en el error $\omega(x) \sim \mathcal{N}\left(0,\frac{|x|}{10}\right)$. 

\begin{figure}[H]
	\centering
	\caption{Datos simulados con heterocedasticidad.}
	\includegraphics[width=0.7\textwidth]{Figures/Simulation/heteroscedasticity/sample.png}
	\label{sample_heteroscedasticity}
\end{figure}

Como se puede observar en la figura \ref{sample_heteroscedasticity}, \'unicamente la mediana de los datos es una l\'inea totalmente recta. Los dem\'as cuantiles se expresan como dos rectas con distinta pendiente, que se unen en el punto $x = 0$. Este comportamiento, similar al valor absoluto, tampoco se puede expresar fielmente con polinomios, sino que s\'olo se puede aproximar. A continuaci\'on se presentan los resultados del ajuste de ambos modelos.

\begin{table}[H]
\centering
\caption{Error cuadrático medio de datos que presentan heterocedasticidad.} 
\begin{tabular}{ccc}
  \hline
Cuantil & Modelo Tradicional & Modelo GPDP \\ 
  \hline
0.95 & 0.51 & 0.58 \\ 
  0.50 & 0.02 & 0.17 \\ 
  0.25 & 0.15 & 0.14 \\ 
   \hline
\end{tabular}
\label{mse_heteroscedasticity}
\end{table}

\begin{table}[H]
\centering
\caption{Correlación al cuadrado de datos que presentan heterocedasticidad.} 
\begin{tabular}{ccc}
  \hline
Cuantil & Modelo Tradicional & Modelo GPDP \\ 
  \hline
0.95 & 0.63 & 0.78 \\ 
  0.50 & 0.98 & 0.86 \\ 
  0.25 & 0.86 & 0.84 \\ 
   \hline
\end{tabular}
\label{corr_heteroscedasticity}
\end{table}

\begin{table}[H]
\centering
\caption{Porcentaje de valores reales dentro del intervalo de confianza de datos que presentan heterocedasticidad.} 
\begin{tabular}{ccc}
  \hline
Cuantil & Modelo Tradicional & Modelo GPDP \\ 
  \hline
0.95 & 68\% & 95\% \\ 
  0.50 & 100\% & 100\% \\ 
  0.25 & 81\% & 100\% \\ 
   \hline
\end{tabular}
\label{within_heteroscedasticity}
\end{table}

\begin{figure}[htbp]
	\centering
	\caption{Ajuste de los modelos Tradicional y \textit{GPDP}, para un conjunto de datos con heterocedasticidad.}
	\includegraphics[width=\textwidth]{Figures/Simulation/heteroscedasticity/predictions.png}
	\captionsetup{singlelinecheck=off,font=footnotesize}
    \caption*{Nota: La l\'inea roja representa el valor real de cada cuantil, la l\'inea azul representa la mediana de la distribuci\'on posterior predictiva y el \'area gris su intervalo de probabilidad al 95\%.}
	\label{models_heteroscedasticity}
\end{figure}

Como se puede observar en la figura \ref{models_heteroscedasticity}, el modelo tradicional interpret\'o el valor de todos los cuantiles como una simple l\'inea recta, \'unicamente variando la ordenada al origen. De hecho, el modelo de la mediana logr\'o una estimaci\'on excelente, como lo reafirman los cuadros \ref{mse_heteroscedasticity}, \ref{corr_heteroscedasticity} y \ref{within_heteroscedasticity}. Sin embargo, conforme la estimaci\'on se fue alejando de la mediana y los valores reales de los cuantiles fueron distingui\'endose de la l\'inea recta, la calidad de la estimaci\'on predictiva se redujo dr\'asticamente.

En cuanto al modelo GPDP, los cuadros \ref{mse_heteroscedasticity}, \ref{corr_heteroscedasticity} y \ref{within_heteroscedasticity} tambi\'en muestran que conforme se tom\'o un cuantil m\'as lejano a la mediana, la calidad de la estimaci\'on empeor\'o. Sin embargo, la ca\'ida fue menos fuerte, particularmente para las medidas de correlaci\'on y porcentaje de valores dentro del intervalo. Esto se debe principalmente a que el modelo tradicional modela la media, y a partir de ah\'i cada cuantil s\'olo representa un cambio de nivel, pero manteniendo la forma de la funci\'on. En cambio, el modelo GPDP contempla que cada cuantil puede tener una forma propia, distinta de los otros cuantiles.

\subsection{Error asim\'etrico}

Hasta este punto se han realizado perturbaciones al modelo tradicional de regresi\'on en la forma de la distribuci\'on de los datos (colas pesadas y heterocedasticidad). Sin embargo, dichas perturbaciones mantuvieron una caracter\'istica del modelo tradicional: la simetr\'ia del error. En el siguiente conjunto de datos se desafi\'o dicha condici\'on, usando un error aleatorio asim\'etrico. Particularmente se utiliz\'o $\omega(x) \sim Gamma(1,1)$.

La funci\'on tendencia involucr\'o a las funciones coseno y exponencial, ambas pudi\'endose aproximar de buena forma con polinomios, utilizando sus series de Taylor. Sin embargo, cabe recordar que \'unicamente se est\'an utilizando polinomios de grado 5 en el modelo tradicional, por lo que podr\'ian resultar insuficientes. Su expresi\'on matem\'atica fue
\begin{equation*}
    g(x) = \frac{1}{5} x cos(x) - \frac{1}{5}exp\left(\frac{x}{10}\right).
\end{equation*}
Los datos obtenidos despu\'es de simular se pueden observar en la figura \ref{sample_asymmetric}.

\begin{figure}[H]
	\centering
	\caption{Datos simulados con error asim\'etrico.}
	\includegraphics[width=0.7\textwidth]{Figures/Simulation/asymmetric/sample.png}
	\label{sample_asymmetric}
\end{figure}

Como se podr\'a dar cuenta el lector, la asimetr\'ia del error provoca cambios en la distancia de los cuantiles respecto a la mediana. Por dar un ejemplo, el cuantil $0.025$\textit{-\'esimo} (el l\'imite inferior del intervalo) est\'a mucho m\'as cerca de la mediana, que el cuantil $0.975$\textit{-\'esimo} (el l\'imite superior del intervalo). Esta situaci\'on es at\'ipica, ya que estamos acostumbrados a que ambos est\'en a la misma distancia, propiciado por la simetr\'ia. Por ello, es pronosticable que el modelo tradicional sufrir\'a para reflejar dicho comportamiento, debido a que supone error normal, el cual presenta la simetr\'ia a la que estamos acostumbrados. A continuaci\'on se presentan los resultados obtenidos del ajuste de ambos modelos.

\begin{table}[H]
\centering
\caption{Error cuadrático medio de datos con error asimétrico.}
\begin{tabular}{ccc}
  \hline
Cuantil & Modelo Tradicional & Modelo GPDP \\ 
  \hline
0.95 & 1.69 & 1.23 \\ 
  0.50 & 1.65 & 0.64 \\ 
  0.25 & 1.53 & 0.44 \\ 
   \hline
\end{tabular}
\label{mse_asymmetric}
\end{table}

\begin{table}[H]
\centering
\caption{Correlación al cuadrado de datos con error asimétrico.}
\begin{tabular}{ccc}
  \hline
Cuantil & Modelo Tradicional & Modelo GPDP \\ 
  \hline
0.95 & 0.00 & 0.50 \\ 
  0.50 & 0.00 & 0.61 \\ 
  0.25 & 0.00 & 0.69 \\ 
   \hline
\end{tabular}
\label{corr_asymmetric}
\end{table}

\begin{table}[H]
\centering
\caption{Porcentaje de valores reales dentro del intervalo de confianza de datos con error asimétrico.} 
\begin{tabular}{ccc}
  \hline
Cuantil & Modelo Tradicional & Modelo GPDP \\ 
  \hline
0.95 & 58\% & 100\% \\ 
  0.50 & 46\% & 100\% \\ 
  0.25 & 52\% & 100\% \\ 
   \hline
\end{tabular}
\label{within_asymmetric}
\end{table}

\begin{figure}[htbp]
	\centering
	\caption{Ajuste de los modelos Tradicional y \textit{GPDP}, para un conjunto de datos con error asim\'etrico.}
	\includegraphics[width=\textwidth]{Figures/Simulation/asymmetric/predictions.png}
	\captionsetup{singlelinecheck=off,font=footnotesize}
    \caption*{Nota: La l\'inea roja representa el valor real de cada cuantil, la l\'inea azul representa la mediana de la distribuci\'on posterior predictiva y el \'area gris su intervalo de probabilidad al 95\%.}
	\label{models_asymmetric}
\end{figure}

La figura \ref{models_asymmetric} permite observar que el modelo tradicional infiere una mediana, en su mayor\'ia, mayor que la real. Esto se debe a que, por el supuesto de simetr\'ia, impl\'icitamente modela a la mediana como si fuera la media, siendo la segunda m\'as grande en su valor real, que la primera. 

La misma figura sugiere que el modelo tradicional interpret\'o a la funci\'on de los diversos cuantiles pr\'acticamente como una constante, y atribuy\'o las variaciones \'unicamente al error aleatorio. La sugerencia de la constante tambi\'en se sustenta en el cuadro \ref{corr_asymmetric}, donde es posible ver que la correlaci\'on es pr\'acticamente 0 (la primer cifra significativa es el tercer decimal), respecto a los valores originales. Adem\'as, la hip\'otesis de que las variaciones se le atribuyeron al error aleatorio se apoya en los bajos porcentajes de valores reales dentro de los intervalos de predicci\'on, presentados en el cuadro \ref{within_asymmetric}.

Por otro lado, se puede verificar en las tablas \ref{mse_asymmetric}, \ref{corr_asymmetric} y \ref{within_asymmetric} que el modelo GPDP logr\'o un mejor desempeño en todas las m\'etricas, para todos los cuantiles. Sin embargo, no hay que perder de vista que de todos los casos vistos hasta este punto, es el que present\'o intervalos de probabilidad m\'as amplios. Es decir, es el caso en el que hay mayor incertidumbre respecto a la predicci\'on. Adem\'as, present\'o niveles de correlaci\'on bajos, comparativamente hablando. No necesariamente la causa fue el error asim\'etrico, sino pudo provenir de que la forma de la funci\'on tendencia es at\'ipica y compleja de inferir para muchos modelos.

\subsection{Discontinuidades}

Dejando atr\'as las modificaciones en el error $\omega(x)$, ahora se presenta una situaci\'on at\'ipica en la funci\'on tendencia. Tanto el modelo tradicional, como el modelo GPDP, suponen continuidad en las funciones de los cuantiles. Por ello, me pareci\'o interesante averiguar c\'omo se comportar\'ia su ajuste, en caso de modelar datos que provinieran de una funci\'on tendencia discontinua. Particularmente se simul\'o utilizando la funci\'on
\begin{equation*}
g(x) = 
\begin{cases}
   - \frac{1}{5} x + 2cos(x) - 2 &\text{ si } x \in (-\infty, -5) \\
   \frac{1}{10} x^2 - \frac{1}{2}x + 1 &\text{ si } x \in [-5, 5) \\
   6 - ln(x) &\text{ si } x \in [-5, \infty)
\end{cases},
\end{equation*}
con el error t\'ipico de los modelos de regresi\'on, $\omega(x) \sim \mathcal{N}(0,1)$.

El comportamiento discontinuo se puede entender de mejor manera gr\'aficamente, como lo muestra la figura \ref{sample_discontinuous}. Adem\'as, en dicha figura tambi\'en se pueden ver los puntos obtenidos de la simulaci\'on.

\begin{figure}[H]
	\centering
	\caption{Datos simulados con discontinuidades.}
	\includegraphics[width=0.7\textwidth]{Figures/Simulation/discontinuous/sample.png}
	\label{sample_discontinuous}
\end{figure}

Como se puede observar en la figura \ref{sample_discontinuous}, y en la expresi\'on matem\'atica antes mencionada, la funci\'on tendencia tiene tres secciones: la primera con un comportamiento oscilatorio, la segunda, un trozo de par\'abola, y la \'ultima, un logaritmo, que por el nivel de la covariable, asemeja a una recta. Todas, por separado, son continuas. Por lo tanto, son aproximables con polinomios en un intervalo cerrado. Sin embargo, entre cada secci\'on hay un salto de discontinuidad, por lo que la funci\'on global no puede ser expresada fielmente usando dicha familia. Ante dicha dificultad se ajustaron los modelos tradicional y GPDP, obteniendo los resultados que se presentan a continuaci\'on. 

\begin{table}[H]
\centering
\caption{Error cuadrático medio de datos simulados con discontinuidades.} 
\begin{tabular}{ccc}
  \hline
Cuantil & Modelo Tradicional & Modelo GPDP \\ 
  \hline
0.95 & 7.08 & 1.10 \\ 
  0.50 & 2.23 & 0.35 \\ 
  0.25 & 3.19 & 0.47 \\ 
   \hline
\end{tabular}
\label{mse_discontinuous}
\end{table}

\begin{table}[H]
\centering
\caption{Correlación al cuadrado de datos simulados con discontinuidades.}
\begin{tabular}{ccc}
  \hline
Cuantil & Modelo Tradicional & Modelo GPDP \\ 
  \hline
0.95 & 0.38 & 0.83 \\ 
  0.50 & 0.38 & 0.90 \\ 
  0.25 & 0.38 & 0.88 \\ 
   \hline
\end{tabular}
\label{corr_discontinuous}
\end{table}

\begin{table}[H]
\centering
\caption{Porcentaje de valores reales dentro del intervalo de confianza de datos simulados con discontinuidades.} 
\begin{tabular}{ccc}
  \hline
Cuantil & Modelo Tradicional & Modelo GPDP \\ 
  \hline
0.95 & 28\% & 98\% \\ 
  0.50 & 46\% & 100\% \\ 
  0.25 & 53\% & 99\% \\ 
   \hline
\end{tabular}
\label{within_discontinuous}
\end{table}

\begin{figure}[htbp]
	\centering
	\caption{Ajuste de los modelos Tradicional y \textit{GPDP}, para un conjunto de datos simulados con discontinuidades.}
	\includegraphics[width=\textwidth]{Figures/Simulation/discontinuous/predictions.png}
	\captionsetup{singlelinecheck=off,font=footnotesize}
    \caption*{Nota: La l\'inea roja representa el valor real de cada cuantil, la l\'inea azul representa la mediana de la distribuci\'on posterior predictiva y el \'area gris su intervalo de probabilidad al 95\%.}
	\label{models_discontinuous}
\end{figure}

Como se puede observar en la figura \ref{models_discontinuous}, el modelo tradicional estuvo muy lejos de inferir correctamente la forma real de la funci\'on de los distintos cuantiles. Esto tambi\'en lo confirma el cuadro \ref{within_discontinuous}, ya que en la mejor de las estimaciones del modelo tradicional, \'unicamente el $53\%$ de los datos reales se pronostic\'o dentro del intervalo de probabilidad.
Al parecer, atribuy\'o el primer salto \'unicamente al error aleatorio, el cual estim\'o de gran magnitud, raz\'on por lo que su estimaci\'on de los cuantiles $0.25$ y $0.95$\textit{-\'esimos} fue m\'as alejada de la mediana de lo que era realmente. El segundo salto parece que s\'i lo detect\'o, pero la elevaci\'on la realiz\'o paulatinamente, llegando al punto m\'aximo despu\'es de 5 unidades de donde realmente ocurri\'o.

Por otro lado, como lo confirma la exploraci\'on visual de la figura \ref{models_discontinuous}, as\'i como los cuadros \ref{mse_discontinuous}, \ref{corr_discontinuous} y \ref{within_discontinuous}, la estimaci\'on del modelo GPDP fue muy superior en calidad, debido a que tuvo errores cuadr\'aticos medios consistentemente bajos, correlaci\'on cuadr\'atica de al menos $83\%$ y al menos $98\%$ de los valores reales adentro del intervalo de probabilidad de la predicci\'on. Debido a la construcci\'on del modelo, tuvo una estimaci\'on continua de la funci\'on de los cuantiles. Sin embargo, lo compens\'o elevando r\'apidamente su nivel (dentro de vecindades m\'as pequeñas que la unidad) cuando hubo saltos de discontinuidad en la funci\'on real. Por ello, los intervalos donde la estimaci\'on fue mala fueron en realidad muy angostos.

\section{Comparaci\'on general entre modelos}
\label{models_comp}

El an\'alisis de los datos anteriores brinda informaci\'on para pensar que el ajuste del modelo GPDP fue, la mayor\'ia de las veces, igual o mejor que el del modelo tradicional de regresi\'on a la media, cuando el objetivo fue estimar el valor de los cuantiles. Esto se dio principalmente cuando se modelaron valores extremos, o los datos provinieron de situaciones que violan los supuestos de la regresi\'on tradicional a la media. 

Pasando a un an\'alisis m\'as minucioso, observ\'e detalles que considero vale la pena resaltar. Para empezar, la magnitud de los intervalos de probabilidad de predicci\'on, ya que los que present\'o el modelo GPDP fueron com\'unmente m\'as anchos que los del tradicional. Esto no significa que su estimaci\'on haya sido peor. De hecho, quiz\'as fue lo contrario. Mientras que los intervalos de probabilidad al $95\%$ del modelo GPDP siempre contuvieron, al menos, al $95\%$ de los valores reales de los cuantiles (como era esperado), hubo intervalos del modelo tradicional al $95\%$ de probabilidad que contuvieron \'unicamente al $9\%$ de los valores reales.

Hablando de los intervalos de probabilidad, es interesante observar que aquellos provenientes del GPDP se hicieron m\'as angostos, conforme hubo m\'as datos observados cercanos a ese valor, y m\'as anchos, ante la ausencia de datos en la vecindad pr\'oxima. Este comportamiento provino de la estructura de correlaci\'on dada por el proceso Gaussiano, recordando que depende de la distancia entre los datos. Por su parte, en el modelo tradicional de regresi\'on a la media, el intervalo de predicci\'on tuvo menor anchura en la parte media de los datos de $x$, y creci\'o muy ligeramente conforme se acerc\'o a los valores m\'inimo y m\'aximo observados. En lo personal, prefiero el comportamiento del modelo GPDP, ya que considero que entre mayor sea la cantidad de datos cercanos, se deber\'ia tener menor incertidumbre, y viceversa.

Tambi\'en es interesante observar que la mediana de la predicci\'on en el modelo tradicional se encontr\'o usualmente en medio del intervalo. Por otro lado, debido a que supone que el error es una mezcla infinita de procesos de Dirichlet, el modelo GPDP present\'o, en algunos casos, asimetr\'ia en la distribuci\'on posterior predictiva.

Las personas que hemos trabajado con el modelo tradicional de regresi\'on a la media estamos acostumbrados a que cualquier estimaci\'on puntual que tomemos es suave, ya que se puede expresar como polinomio, y, por lo tanto, es infinitamente diferenciable. En cambio, debido a los supuestos que introduce el modelo GPDP para el proceso Gaussiano, sus estimaciones puntuales mantuvieron un comportamiento mucho menos suave, pero que fue \'util para capturar cambios bruscos en la funci\'on. Sin embargo, es importante recordar que la dificultad de interpretaci\'on que tienen los modelos no param\'etricos representa una de sus mayores desventajas.

Finalmente, es interesante comparar la diferencia en tiempos de ajuste y predicci\'on para ambos modelos, los cuales se muestran en las figuras \ref{fit_time} y \ref{pred_time}. \footnote{Estos tiempos se obtuvieron corriendo los procesos en una MacBook, con procesador 1.1 GHz Intel Core M, y memoria RAM DDR3, de 8 GB y 1600 MHz. Los procesos de predicci\'on se corrieron en paralelo, usando 2 de los 4 núcleos.} Cabe señalar que los tiempos presentados son el acumulado de ajustar o predecir los modelos de los tres cuantiles.

\begin{table}[H]
\centering
\caption{Tiempo de ajuste por conjunto de datos, para cada modelo.} 
\begin{tabular}{ccc}
  \hline
Datos & Tradicional (seg) & GPDP (seg) \\ 
  \hline
Supuestos tradicionales & menos de 1 & 2,498 \\ 
  Colas pesadas & menos de 1 & 4,006 \\ 
  Heterocedasticidad & menos de 1 & 3,502 \\ 
  Error asimétrico & menos de 1 & 6,707 \\ 
  Discontinuidades & menos de 1 & 3,062 \\ 
   \hline
\end{tabular}
\label{fit_time}
\end{table}

\begin{table}[H]
\centering
\caption{Tiempo de predicción por conjunto de datos, para cada modelo.} 
\begin{tabular}{ccc}
  \hline
Datos & Tradicional (seg) & GPDP (seg) \\ 
  \hline
Supuestos tradicionales & 6.4 & 563.8 \\ 
  Colas pesadas & 4.7 & 529.0 \\ 
  Heterocedasticidad & 5.3 & 534.5 \\ 
  Error asimétrico & 5.7 & 537.4 \\ 
  Discontinuidades & 5.0 & 533.2 \\ 
   \hline
\end{tabular}
\label{pred_time}
\end{table}

Analizando los tiempos de ajuste, y recordando que se usaron 60 datos en todos los casos, es clara la bondad dada por la propiedad conjugada del modelo tradicional. \'Esta permite obtener una expresi\'on anal\'itica cerrada para la distribuci\'on posterior de los par\'ametros, y por ende, el ajuste consiste \'unicamente en multiplicaciones matriciales que toman menos de un segundo. Adem\'as, por las caracter\'isticas de este modelo, \'unicamente se requiri\'o ajustar una vez por conjunto de datos, independientemente de que se usara para varios cuantiles.

Por otro lado, el modelo GPDP requiere un ajuste distinto por cuantil, y el tiempo que tarda en correr est\'a siempre en funci\'on de la cantidad de iteraciones que se eligen realizar. En el caso particular de este trabajo, se corrieron 20,000 iteraciones para todos los casos, lo cual representa un n\'umero poco sobresaliente, comparado con las que normalmente se utilizan cuando se busca tener una precisi\'on alta. El modelo que menos tard\'o requiri\'o m\'as de 40 minutos para realizar el ajuste, por lo que la diferencia es abismal entre ambos modelos.

De hecho, es interesante ver c\'omo hay grandes diferencias dentro de los propios ajustes del modelo GPDP. La raz\'on principal es la simulaci\'on de una distribuci\'on normal truncada en cada iteraci\'on del algoritmo, usando un m\'etodo de aceptaci\'on y rechazo. Cuando los datos cumplieron los supuestos tradicionales, fue cuando le cost\'o menos trabajo realizar las simulaciones. En cambio, cuando se us\'o una funci\'on con muchas subidas y bajadas, y error asim\'etrico, el algoritmo requiri\'o m\'as del doble de tiempo para poder ajustar el modelo.

Para el caso de las predicciones, el n\'umero de puntos aument\'o a 300 para todos los casos. Sin embargo, debido a que estas simulaciones ya no pertenecen a una cadena de Markov, se pueden correr en paralelo para los dos modelos. Ambos requieren simular una distribuci\'on Normal, pero el tradicional simula una marginal por cada dato, mientras que el GPDP requiere simular una conjunta, de dimensi\'on igual al n\'umero de datos. Esa es la raz\'on por la que tarda 100 veces m\'as.

As\'i, mientras que el modelo GPDP mostr\'o, en general, mejores resultados en la precisi\'on, el modelo tradicional de regresi\'on a la media fue mucho m\'as efectivo en tiempo. Por ello, el modelador deber\'ia considerar cu\'al de ambas prefiere en su estudio, para poder tomar una mejor decisi\'on de qu\'e modelo elegir.


\newpage
\chapter[Conclusiones y trabajo futuro]{Conclusiones y trabajo futuro}

Si bien los modelos de regresi\'on a la media han sido de mucha utilidad en las \'ultimas d\'ecadas, principalmente cuando el poder computacional era menor, es importante darse cuenta que actualmente existen contextos en los que resultan insuficientes. Ya sea porque sea desea la mejor aproximaci\'on posible de alg\'un estad\'istico distinto a la media, o debido a que hay ocasiones en las que no se cumplen algunos de sus supuestos.

De manera similar, la relaci\'on lineal en los par\'ametros y la distribuci\'on Normal del error han sido fundamentales para que los modelos de regresi\'on hayan proliferado en una gran cantidad de industrias, tanto por su interpretabilidad, como por su bajo costo de estimaci\'on. Sin embargo, es imposible ignorar que \'unicamente representan un subconjunto del universo de funciones y errores aleatorios posibles. Crear modelos que permitan una mayor flexibilidad, como aquellos que utilizan m\'etodos no param\'etricos, lograr\'a una representaci\'on m\'as certera de la realidad de la que provienen los datos.

Al momento de hacer aproximaciones estad\'isticas, se dice que utilizar el paradigma Bayesiano muchas veces presenta la ventaja de poder introducir el conocimiento previo de las personas expertas en el fen\'omeno a estudiar. Desafortunadamente, en un modelo jer\'arquico, como lo es el expuesto en esta tesis, es complicado transmitir dicho conocimiento. Principalmente, debido a que los hiper-par\'ametros se encuentran varias capas abajo de aquellos par\'ametros que tienen una interpretaci\'on natural. A pesar de ello, vale la pena resaltar que dicha construcci\'on jer\'arquica, con la flexibilidad que brinda, es posible y es congruente gracias al paradigma Bayesiano, debido a que todas las expresiones son completamente probabil\'isticas y fundadas en un cuerpo axiom\'atico. 

Un reto importante que present\'o este trabajo fue el desarrollo del paquete en R para implementar el modelo GPDP. Primero, porque se requiri\'o plantear te\'oricamente la distribuciones condicionales necesarias para que corriera el simulador de Gibbs. Y segundo, porque tuve que buscar programar de forma general y eficiente, para que el paquete funcionara siempre que recibiera los par\'ametros predefinidos, y corriera lo m\'as r\'apido posible, ante la desventaja que representa s\'olo poder calcular a la vez una iteraci\'on de la cadena de Markov. De hecho, dej\'e el n\'umero de iteraciones como un par\'ametro a elecci\'on del modelador, para que pueda decidir si prefiere precisi\'on o velocidad.

Si bien estos avances son significativos, a\'un existe mucho que explorar respecto a lo expuesto en esta tesis. Por ejemplo, el modelo planteado en este trabajo no es capaz de darle un peso distinto a cada variable explicativa, sino que las toma  a todas por igual al momento de calcular la distancia entre observaciones. Para mejorar esta situaci\'on se podr\'ia plantear una descomposici\'on de la funci\'on del cuantil en la suma de varios procesos Gaussianos, uno por covariable, lo que brindar\'ia un mayor peso a aquellas que en efecto sean m\'as significativas para explicar el fen\'omeno en cuesti\'on. 

Adem\'as, ser\'ia conveniente la inclusi\'on de un par\'ametro de rango que regule din\'amicamente la relaci\'on entre la distancia y la covarianza entre observaciones. Por ejemplo, a\'un cuando est\'en estandarizados los datos, una misma distancia podr\'ia significar una covarianza grande entre observaciones para alguna covariable o fen\'omeno, pero covarianza casi nula para otro. Lograr implementar este par\'ametro din\'amico seguramente mejorar\'a el ajuste.

Finalmente, tendr\'ia una gran utilidad el desarrollar una medida robusta de bondad de ajuste para este tipo de modelos. Esto brindar\'ia cualidades importantes al modelo GPDP, como el poder hacer selecci\'on de variables, y tambi\'en permitir\'ia saber qu\'e tan bueno o malo es el modelo, en comparaci\'on con lo dem\'as disponibles. 

\newpage 

%%%%%%%%%%%%%%%%%%%%%%%%%%%%%%%%%%%%%%%%%%%%%%%%%%%%%%%%%%%%%
%%    Bibliografia
%%%%%%%%%%%%%%%%%%%%%%%%%%%%%%%%%%%%%%%%%%%%%%%%%%%%%%%%%%%%%

\nocite{*} %Even non-cited BibTeX-Entries will be shown.
\bibliographystyle{authordate1}
%Style of Bibliography: plain / apalike / amsalpha / ...
\bibliography{Bibliography} 

%%%%%%%%%%%%%%%%%%%%%%%%%%%%%%%%%%%%%%%%%%%%%%%%%%%%%%%%%%%%%
%%  Apendices
%%%%%%%%%%%%%%%%%%%%%%%%%%%%%%%%%%%%%%%%%%%%%%%%%%%%%%%%%%%%%
\appendix
\chapter[Distribuciones de probabilidad]{Distribuciones de probabilidad }\label{chap:Distributions}

\section{Distribuci\'on Normal condicional}

\begin{prop*}
    Sea $X \in \mathbb{R}^m$ un vector aleatorio que tiene distribuci\'on Normal conjunta y est\'a particionado de la siguiente manera:

    \begin{equation*}
    \begin{aligned}
        X = 
        &\left[
        \begin{array}{c}
            X_1  \\
            X_2
        \end{array}
        \right], \\
        \text{ con dimensiones }
        &\left[
        \begin{array}{c}
            (m-q)  \\
            q
        \end{array}
        \right].
    \end{aligned}
    \end{equation*}
    
    Entonces, la media $\mu \in \mathbb{R}^m$ y varianza $\Sigma \in \mathbb{R}^{m \times m}$ de $X$ se pueden escribir
    \begin{equation*}
    \begin{aligned}
        \mu = 
        &\left[
        \begin{array}{c}
            \mu_1  \\
            \mu_2
        \end{array}
        \right], \\
        \text{ con dimensiones }
        &\left[
        \begin{array}{c}
            (m-q)  \\
            q
        \end{array}
        \right], y\\
        \Sigma = 
        &\left[
        \begin{array}{cc}
            \Sigma_{11} & \Sigma_{12}  \\
            \Sigma_{21} & \Sigma_{22}
        \end{array}
        \right], \\
        \text{ con dimensiones }
        &\left[
        \begin{array}{cc}
            (m-q) \times (m-q)  & (m-q) \times q  \\
            q \times (m-q) & q \times q
        \end{array}
        \right].
    \end{aligned}
    \end{equation*}
    
    La distribuci\'on condicional de $X_2$, sujeta a que $X_1 = a$ es Normal con $X_2|X_1=a \sim \mathcal{N}(X_2|\bar{\mu},\bar{\Sigma})$, donde
    
    \begin{equation*}
    \begin{aligned}
        \bar{\mu} &= \mu_2 + \Sigma_{2,1}\Sigma_{11}^{-1}(a-\mu_1) \\
        \bar{\Sigma} &= \Sigma_{22} - \Sigma_{21}\Sigma_{11}^{-1}\Sigma_{12}.
    \end{aligned}
    \end{equation*}
\end{prop*}

\section{Distribuci\'on de Dirichlet}

\begin{defin}
    Se dice que un vector aleatorio $x \in \mathbb{R}^n$ se distribuye de acuerdo a la \textbf{distribuci\'on de Dirichlet}  $\mathbf{(x \sim Dir(\alpha))}$ con vector de par\'ametros $\alpha$, espec\'ificamente,
    \begin{equation*}
        x = 
        \left(\begin{array}{c}
            x_1  \\
            \vdots \\
            x_n
        \end{array}\right),
        \qquad
        \alpha = 
        \left(\begin{array}{c}
            \alpha_1  \\
            \vdots \\
            \alpha_n
        \end{array}\right),
    \end{equation*}
    para los cuales se cumplen las restricciones
    \begin{equation*}
    \begin{aligned}
        x_i > 0, \forall i &\in \{1,...,n\} \\
        \sum_{i=1}^n x_i &= 1 \\
        \alpha_i > 0, \forall i &\in \{1,...,n\},
    \end{aligned}
    \end{equation*}
    si su funci\'on de densidad  es
    \begin{equation*}
        f(x|\alpha) = 
        \frac {1}{\mathrm {B} (\alpha)}
        \prod _{i=1}^{n}x_{i}^{\alpha _{i}-1},
    \end{equation*}
    donde $\mathrm{B}$ es la funci\'on Beta multivariada, y puede ser expresada en t\'erminos de la funci\'on $\Gamma$ como 
    \begin{equation*}
       \mathrm{B}(\alpha)=
       \frac {\prod _{i=1}^{n}\Gamma (\alpha _{i})}
       {\Gamma \left(\sum _{i=1}^{n}\alpha _{i}\right)},
       \qquad 
       \alpha =(\alpha _1,\cdots ,\alpha _n). 
    \end{equation*}
    
    La esperanza y varianza de cada $x_i$ son los siguientes:
    
    \begin{equation*}
    \begin{aligned}
        \mathbb{E}[x_i] &= \frac{\alpha_i}{\sum_{k=1}^n \alpha_k} \\
        Var(x_i) &= \frac
        {\alpha_i \left( \sum_{k=1}^n \alpha_k - \alpha_i \right)}
        {\left( \sum_{k=1}^n \alpha_k \right)^2 \left( \sum_{k=1}^n \alpha_k + 1 \right)}
    \end{aligned}
    \end{equation*}
    
\end{defin}

Es com\'un que esta distribuci\'on sea usada como la inicial conjugada de la distribuci\'on multinomial, debido a que el vector $x$ tiene las mismas propiedades de una distribuci\'on de probabilidad discreta (elementos positivos y que en conjunto suman 1).
\input{Appendices/B-MCMC}

\end{document}

